\chapter{A integral de Riemann} \label{int:chapter}

%%%%%%%%%%%%%%%%%%%%%%%%%%%%%%%%%%%%%%%%%%%%%%%%%%%%%%%%%%%%%%%%%%%%%%%%%%%%%%

\section{A integral de Riemann}
\label{sec:rint}

%mbxINTROSUBSECTION

\sectionnotes{1,5 aulas}

Uma integral é uma maneira de \myquote{somar} os valores de uma função.
Às vezes, estudantes de cálculo confundem a \emph{integral} com a
\emph{primitiva}.
Informalmente, a integral é a área sob a curva, e nada mais.
O fato de podermos calcular uma primitiva usando a integral é um resultado
nada trivial, que precisamos demonstrar.
Vamos definir a \emph{integral de Riemann}%
\footnote{Nomeada em homenagem ao matemático alemão
\href{https://en.wikipedia.org/wiki/Riemann}{Georg Friedrich Bernhard Riemann}
(1826--1866).}
usando a integral de Darboux%
\footnote{Nomeada em homenagem ao matemático francês
\href{https://en.wikipedia.org/wiki/Darboux}{Jean-Gaston Darboux} (1842--1917).},
uma definição equivalente, mas tecnicamente mais simples.

\subsection{Partições e integrais inferior e superior}

Queremos integrar uma função limitada definida em um intervalo $[a,b]$.
Primeiro definimos duas integrais auxiliares, que estão definidas para todas
as funções limitadas. Só então podemos falar da integral de Riemann e das
funções que ela pode integrar, isto é, das funções integráveis à Riemann.

\begin{defn}
Uma \emph{\myindex{partição}} $P$ de $[a,b]$ é
um conjunto finito de números $\{ x_0,x_1,x_2,\ldots,x_n \}$ tal que
\begin{equation*}
a = x_0 < x_1 < x_2 < \cdots < x_{n-1} < x_n = b .
\end{equation*}
Escrevemos
\begin{equation*}
\Delta x_i \coloneqq x_i - x_{i-1} .
\end{equation*}
Suponha que $f \colon [a,b] \to \R$ seja limitada e que $P$ seja uma partição
de $[a,b]$.
Definimos
\begin{align*}
& m_i \coloneqq \inf \, \bigl\{ f(x) : x_{i-1} \leq x \leq x_i \bigr\} , &
& M_i \coloneqq \sup \, \bigl\{ f(x) : x_{i-1} \leq x \leq x_i \bigr\} , \\
& L(P,f) \coloneqq
\sum_{i=1}^n m_i \Delta x_i , &
& U(P,f) \coloneqq
\sum_{i=1}^n M_i \Delta x_i .
\avoidbreak
\end{align*}
Chamamos $L(P,f)$ de \emph{\myindex{soma inferior de Darboux}} e
$U(P,f)$ de \emph{\myindex{soma superior de Darboux}}\index{soma de Darboux}.
\glsadd{not:lowerdarbouxsum}
\glsadd{not:upperdarbouxsum}
\end{defn}

A ideia geométrica das somas de Darboux está indicada na
\figureref{darbouxfig}. A soma inferior é a área dos retângulos
sombreados, e a soma superior é a área de todos os retângulos, incluindo
as partes sombreadas e as não sombreadas. A largura do $i$-ésimo retângulo é
$\Delta x_i$, a altura do retângulo sombreado é $m_i$ e a altura do retângulo
inteiro é $M_i$.

\begin{myfigureht}
\myincludegraphics{darbouxfig}{%
Diagrama do gráfico de uma função, representado por uma linha preta espessa
inteiramente acima do eixo $x$.
O domínio está dividido em 8 intervalos, cujas extremidades estão marcadas,
da esquerda para a direita, de $x_0$ a $x_8$.
O intervalo entre $x_4$ e $x_5$ está indicado como tendo largura
$\Delta x_5$.
Acima de cada intervalo há dois retângulos. Por exemplo,
há um retângulo sombreado de largura $\Delta x_5$,
cujo lado inferior coincide exatamente com o intervalo de
$x_4$ a $x_5$ e cuja altura é $m_5$, o valor mínimo da função nesse
intervalo. O retângulo se prolonga para cima, sem sombreamento, até
$M_5$, o valor máximo da função no intervalo.}
\caption{Exemplos de somas de Darboux.\label{darbouxfig}}
\end{myfigureht}

\begin{prop} \label{sumulbound:prop}
Seja $f \colon [a,b] \to \R$ uma função limitada. Sejam $m, M \in \R$
tais que $m \leq f(x) \leq M$ para todo $x \in [a,b]$. Então, para toda
partição $P$ de $[a,b]$,
\begin{equation}
\label{sumulbound:eq}
m(b-a) \leq
L(P,f) \leq U(P,f)
\leq M(b-a) .
\end{equation}
\end{prop}

\begin{proof}
Seja $P$ uma partição de $[a,b]$. Observe que, para todo $i$,
temos $m \leq m_i \leq M_i \leq M$. Além disso,
$\sum_{i=1}^n \Delta x_i = (b-a)$. Portanto,
\begin{multline*}
m(b-a) =
m \left( \sum_{i=1}^n \Delta x_i \right)
=
\sum_{i=1}^n m \Delta x_i
\leq
\sum_{i=1}^n m_i \Delta x_i 
\\
\leq
\sum_{i=1}^n M_i \Delta x_i
\leq
\sum_{i=1}^n M \Delta x_i 
=
M \left( \sum_{i=1}^n \Delta x_i \right)
=
M(b-a) .
\end{multline*}
Logo, obtemos \eqref{sumulbound:eq}. Em particular, os conjuntos das somas
inferiores e superiores são limitados.
\end{proof}

\begin{defn}
Como os conjuntos das somas inferiores e superiores de Darboux são limitados,
definimos
\begin{align*}
& \underline{\int_a^b} f(x)\,dx \coloneqq
\sup \, \bigl\{ L(P,f) : P \text{ é uma partição de } [a,b] \bigr\} , \\
& \overline{\int_a^b} f(x)\,dx \coloneqq
\inf \, \bigl\{ U(P,f) : P \text{ é uma partição de } [a,b] \bigr\} .
\end{align*}
Chamamos $\underline{\int}$\glsadd{not:lowerdarboux}
de \emph{\myindex{integral inferior de Darboux}} e
$\overline{\int}$\glsadd{not:upperdarboux} de
\emph{\myindex{integral superior de Darboux}}\index{integral de Darboux}.
Para não nos preocuparmos com a variável de integração,
frequentemente escrevemos apenas
\begin{equation*}
\underline{\int_a^b} f \coloneqq
\underline{\int_a^b} f(x)\,dx 
\qquad \text{e} \qquad
\overline{\int_a^b} f \coloneqq
\overline{\int_a^b} f(x)\,dx  .
\end{equation*}
\end{defn}

Para que a integração faça sentido, as integrais inferior e superior de
Darboux deveriam ser o mesmo número, pois queremos ter um único número que
possamos chamar de \emph{integral}. No entanto, essas duas integrais podem ser
diferentes para algumas funções.

\begin{example} \label{example:dirichletfunc}
Considere a \myindex{função de Dirichlet}
$f \colon [0,1] \to \R$, definida por $f(x) \coloneqq 1$ se
$x \in \Q$ e $f(x) \coloneqq 0$ se $x \notin \Q$. Então,
\begin{equation*}
\underline{\int_0^1} f = 0 \qquad \text{e} \qquad
\overline{\int_0^1} f = 1 .
\end{equation*}
Isso ocorre porque, para qualquer partição $P$ e todo $i$, temos
$m_i = \inf \bigl\{ f(x) : x \in [x_{i-1},x_i] \bigr\} = 0$ e
$M_i = \sup \bigl\{ f(x) : x \in [x_{i-1},x_i] \bigr\} = 1$. Logo,
\begin{equation*}
L(P,f) = \sum_{i=1}^n 0 \cdot \Delta x_i = 0 , \quad \text{e} \quad
U(P,f) = \sum_{i=1}^n 1 \cdot \Delta x_i = \sum_{i=1}^n \Delta x_i = 1  .
\end{equation*}
\end{example}

\begin{remark}
Usamos a mesma definição de $\underline{\int_a^b} f$ e
$\overline{\int_a^b} f$
quando $f$ está definida em um conjunto maior $S$ tal que
$[a,b] \subset S$. Nesse caso, usamos a restrição de $f$ a $[a,b]$
e devemos garantir que essa restrição seja limitada em $[a,b]$.
\end{remark}

Para calcular a integral, frequentemente tomamos uma partição $P$ e a
refinamos. Isto é, subdividimos os intervalos da partição em partes ainda
menores.

\begin{defn}
Sejam $P = \{ x_0, x_1, \ldots, x_n \}$ e
$\widetilde{P} = \{ \widetilde{x}_0, \widetilde{x}_1, \ldots,
\widetilde{x}_{\ell} \}$ sejam
partições de $[a,b]$. Dizemos que $\widetilde{P}$ é um
\emph{refinamento}\index{refinamento de uma partição} de $P$
se, como conjuntos, $P \subset \widetilde{P}$.
\end{defn}

Isto é, $\widetilde{P}$ é um refinamento de uma partição quando contém todos
os números de $P$ e, possivelmente, outros números entre eles. Por exemplo,
$\{ 0, 0.5, 1, 2 \}$ é uma partição de $[0,2]$, e
$\{ 0, 0.2, 0.5, 1, 1.5, 1.75, 2 \}$ é um refinamento.
A principal razão para introduzir refinamentos é a seguinte proposição.

\begin{prop} \label{prop:refinement}
Seja $f \colon [a,b] \to \R$ uma função limitada, e seja $P$
uma partição de $[a,b]$. Seja $\widetilde{P}$ um refinamento de $P$.
Então
\begin{equation*}
L(P,f) \leq L(\widetilde{P},f) 
\qquad \text{e} \qquad
U(\widetilde{P},f) \leq U(P,f) .
\end{equation*}
\end{prop}

\begin{proof}
O ponto delicado desta demonstração é estabelecer corretamente a notação.
Seja $\widetilde{P} = \{ \widetilde{x}_0, \widetilde{x}_1, \ldots,
\widetilde{x}_{\ell} \}$ um refinamento de
$P = \{ x_0, x_1, \ldots, x_n \}$. Então
$x_0 = \widetilde{x}_0$ e
$x_n = \widetilde{x}_{\ell}$. De fato, existem inteiros
$k_0 < k_1 < \cdots < k_n$ tais que $x_i = \widetilde{x}_{k_i}$ para
$i=0,1,2,\ldots,n$.

Defina $\Delta \widetilde{x}_q \coloneqq \widetilde{x}_q - \widetilde{x}_{q-1}$ para
$q=0,1,2,\ldots,\ell$.
Veja a \figureref{fig:refinement}.
Obtemos
\begin{equation*}
\Delta x_i
=
x_i - x_{i-1} =
\widetilde{x}_{k_i} - \widetilde{x}_{k_{i-1}} =
\sum_{q=k_{i-1}+1}^{k_i} 
\widetilde{x}_{q} - \widetilde{x}_{q-1}
=
\sum_{q=k_{i-1}+1}^{k_i} \Delta \widetilde{x}_q .
\end{equation*}
\begin{myfigureht}
\myincludepdft{figrefinement}{%
A figura mostra um intervalo de extremidades $x_{i-1}$ e $x_i$ e largura
$\Delta x_i$.
Ele está dividido em três intervalos menores.
A extremidade esquerda $x_{i-1}$
também está indicada por $\widetilde{x}_{q-3}$
e por $\widetilde{x}_{k_{i-1}}$.
De modo semelhante, a extremidade direita $x_i$
também está indicada por $\widetilde{x}_q$ e por $\widetilde{x}_{k_i}$.
Os dois pontos interiores estão indicados por
$\widetilde{x}_{q-2}$ e
$\widetilde{x}_{q-1}$.
Os três intervalos têm comprimentos
$\Delta \widetilde{x}_{q-2}$,
$\Delta \widetilde{x}_{q-1}$
e
$\Delta \widetilde{x}_q$.}
\caption{Refinamento de um subintervalo. Observe que $\Delta x_i =
\Delta \widetilde{x}_{q-2} +
\Delta \widetilde{x}_{q-1} +
\Delta \widetilde{x}_{q}$,
e também que
$k_{i-1}+1 = q-2$ e
$k_{i} = q$.\label{fig:refinement}}
\end{myfigureht}

Seja $m_i$ como antes, correspondente à partição $P$.
Defina $\widetilde{m}_q \coloneqq \inf \bigl\{ f(x) : \widetilde{x}_{q-1} \leq x \leq
\widetilde{x}_q \bigr\}$.
Temos $m_i \leq \widetilde{m}_q$ para $k_{i-1} < q \leq k_i$. Portanto,
\begin{equation*}
m_i \Delta x_i
=
m_i \sum_{q=k_{i-1}+1}^{k_i} \Delta \widetilde{x}_q
=
\sum_{q=k_{i-1}+1}^{k_i} m_i \Delta \widetilde{x}_q
\leq
\sum_{q=k_{i-1}+1}^{k_i} \widetilde{m}_q \Delta \widetilde{x}_q .
\end{equation*}
Assim,
\begin{equation*}
L(P,f) =
\sum_{i=1}^n m_i \Delta x_i
\leq
\sum_{i=1}^n \,
\sum_{q=k_{i-1}+1}^{k_i} \widetilde{m}_q \Delta \widetilde{x}_q
=
\sum_{q=1}^{\ell}
\widetilde{m}_q \Delta \widetilde{x}_q = L(\widetilde{P},f).
\end{equation*}

A demonstração de $U(\widetilde{P},f) \leq U(P,f)$ fica como exercício.
\end{proof}

Com os refinamentos à disposição, provamos o seguinte.
O ponto principal desta proposição é que a integral inferior de Darboux é
menor ou igual à integral superior de Darboux.

\begin{prop} \label{intulbound:prop}
Seja $f \colon [a,b] \to \R$ uma função limitada. Sejam $m, M \in \R$
tais que $m \leq f(x) \leq M$ para todo $x \in [a,b]$. Então
\begin{equation}
\label{intulbound:eq}
m(b-a) \leq
\underline{\int_a^b} f \leq \overline{\int_a^b} f
\leq M(b-a) .
\end{equation}
\end{prop}

\begin{proof}
Pela \propref{sumulbound:prop}, para toda partição $P$,
\begin{equation*}
m(b-a) \leq L(P,f) \leq U(P,f) \leq M(b-a).
\end{equation*}
Da desigualdade
$m(b-a) \leq L(P,f)$ implies $m(b-a) \leq \underline{\int_a^b} f$.
Da desigualdade
$U(P,f) \leq M(b-a)$ implies $\overline{\int_a^b} f \leq M(b-a)$.

A desigualdade central em
\eqref{intulbound:eq} é o ponto principal da proposição.
Sejam $P_1, P_2$ partições de $[a,b]$. Defina
$\widetilde{P} \coloneqq P_1 \cup P_2$.
O conjunto $\widetilde{P}$ é uma partição de $[a,b]$ e um refinamento
tanto de $P_1$ quanto de $P_2$.
Pela \propref{prop:refinement},
$L(P_1,f) \leq L(\widetilde{P},f)$ e
$U(\widetilde{P},f) \leq U(P_2,f)$. Assim,
\begin{equation*}
L(P_1,f) \leq L(\widetilde{P},f) \leq U(\widetilde{P},f) \leq U(P_2,f) .
\end{equation*}
Em outras palavras, para duas partições arbitrárias $P_1$ e $P_2$, temos
$L(P_1,f) \leq U(P_2,f)$.
Lembrando a \propref{infsupineq:prop}, tomamos o supremo e o ínfimo sobre
todas as partições:
\begin{multline*}
\underline{\int_a^b} f = 
\sup \, \bigl\{ L(P,f) : P \text{ é uma partição de } [a,b] \bigr\}
\\
\leq
\inf \, \bigl\{ U(P,f) : P \text{ é uma partição de } [a,b] \bigr\}
=
\overline{\int_a^b} f . \qedhere
\end{multline*}
\end{proof}

\subsection{Integral de Riemann}

Finalmente podemos definir a integral de Riemann. No entanto, ela só está
definida para certa classe de funções, chamadas funções integráveis à
Riemann.

\begin{defn}
Seja $f \colon [a,b] \to \R$ uma função limitada tal que
\begin{equation*}
\underline{\int_a^b} f(x)\,dx = \overline{\int_a^b} f(x)\,dx .
\end{equation*}
Então dizemos que $f$ é \emph{\myindex{integrável à Riemann}}.
O conjunto das funções integráveis à Riemann em $[a,b]$ é denotado
por $\sR\bigl([a,b]\bigr)$.\glsadd{not:integrablefunc}
Quando $f \in \sR\bigl([a,b]\bigr)$, definimos\glsadd{not:riemannint}
\begin{equation*}
\int_a^b f(x)\,dx \coloneqq 
\underline{\int_a^b} f(x)\,dx = \overline{\int_a^b} f(x)\,dx .
\end{equation*}
Como antes, frequentemente escrevemos
\begin{equation*}
\int_a^b f \coloneqq \int_a^b f(x)\,dx.
\end{equation*}
O número $\int_a^b f$ é chamado de \emph{\myindex{integral de Riemann}}
de $f$ ou, às vezes, simplesmente de \emph{integral} de $f$.
\end{defn}

Por definição, uma função integrável à Riemann é limitada.
Recorrendo à \propref{intulbound:prop}, obtemos imediatamente
a proposição a seguir. Veja também a \figureref{fig:integralminmax}.

\begin{prop} \label{intbound:prop}
Seja $f \colon [a,b] \to \R$ uma função integrável à Riemann.
Sejam $m, M \in \R$
tais que $m \leq f(x) \leq M$ para todo $x \in [a,b]$. Então
\begin{equation*}
m(b-a) \leq
\int_a^b f
\leq M(b-a) .
\end{equation*}
\end{prop}
\begin{myfigureht}
\myincludegraphics{integralminmax}{%
Diagrama de um retângulo no plano,
com extensão horizontal entre $a$ e $b$ e extensão vertical entre o eixo $x$
e o nível $M$. O retângulo está dividido em uma parte superior branca,
que contém o gráfico de uma função, e uma parte inferior sombreada.
A reta horizontal que as separa está marcada com $m$.}
\caption{A área sob a curva é limitada superiormente pela área do retângulo
inteiro, $M(b-a)$, e inferiormente pela área da parte sombreada,
$m(b-a)$.\label{fig:integralminmax}}
\end{myfigureht}

Uma forma mais fraca desta proposição é frequentemente útil:
Se $\babs{f(x)} \leq M$ para todo $x \in [a,b]$, então
\begin{equation*}
\abs{\int_a^b f} \leq M(b-a) .
\end{equation*}

\begin{example}
Integramos funções constantes usando a
\propref{intulbound:prop}.
Se $f(x) \coloneqq c$ para alguma constante $c$, tomamos $m = M = c$.
Na desigualdade \eqref{intulbound:eq},
todas as desigualdades devem ser igualdades.
Logo, $f$ é integrável em $[a,b]$ e $\int_a^b f = c(b-a)$.
\end{example}

\begin{example}
Seja $f \colon [0,2] \to \R$ definida por
\begin{equation*}
f(x) \coloneqq
\begin{cases}
1               & \text{se } x < 1,\\
\nicefrac{1}{2} & \text{se } x = 1,\\
0               & \text{se } x > 1.
\end{cases}
\end{equation*}
Afirmamos que $f$ é integrável à Riemann e que $\int_0^2 f = 1$.

Demonstração: Seja $0 < \epsilon < 1$ arbitrário.
Seja $P \coloneqq \{0, 1-\epsilon, 1+\epsilon, 2\}$ uma partição. Usamos a
notação da definição das somas de Darboux. Então
\begin{align*}
m_1 &= \inf  \bigl\{ f(x) : x \in [0,1-\epsilon] \bigr\} = 1 , & 
M_1 &= \sup  \bigl\{ f(x) : x \in [0,1-\epsilon] \bigr\} = 1 , \\
m_2 &= \inf  \bigl\{ f(x) : x \in [1-\epsilon,1+\epsilon] \bigr\} = 0 , & 
M_2 &= \sup  \bigl\{ f(x) : x \in [1-\epsilon,1+\epsilon] \bigr\} = 1 , \\
m_3 &= \inf  \bigl\{ f(x) : x \in [1+\epsilon,2] \bigr\} = 0 , & 
M_3 &= \sup  \bigl\{ f(x) : x \in [1+\epsilon,2] \bigr\} = 0 .
\end{align*}
Além disso, $\Delta x_1 = 1-\epsilon$, $\Delta x_2 = 2\epsilon$ e
$\Delta x_3 = 1-\epsilon$.
Veja a \figureref{darbouxfigstep}.
\begin{myfigureht}
\myincludegraphics{darbouxfigstep}{%
Diagrama do gráfico de uma função que vale 1 de 0 a 1,
vale um meio em 1 e vale 0 de 1 a 2.
O eixo horizontal está dividido em três intervalos:
de 0 a 1 menos épsilon, de 1 menos épsilon a 1 mais épsilon e de 1 mais
épsilon a 2.
Seus comprimentos são $\Delta x_1 = 1-\epsilon$,
$\Delta x_2 = 2\epsilon$ e
$\Delta x_3 = 1-\epsilon$, respectivamente.
Há um retângulo sombreado de altura 1 para $x$ entre 0 e
1 menos épsilon. Há um retângulo branco de altura 1
para $x$ entre 1 menos épsilon e 1 mais épsilon.
No eixo vertical,
o ponto inferior está marcado como $M_3=m_2=m_3=0$.
O ponto superior está marcado como $M_1=M_2=m_1=1$.}
\caption{Somas de Darboux para a função em degraus. $L(P,f)$ é a área do
retângulo sombreado, $U(P,f)$ é a área dos dois retângulos, e
$U(P,f)-L(P,f)$ é a área do retângulo não sombreado.\label{darbouxfigstep}}
\end{myfigureht}

Calculamos
\begin{align*}
& L(P,f) = \sum_{i=1}^3 m_i \Delta x_i =
1 \cdot (1-\epsilon) + 0 \cdot 2\epsilon + 0 \cdot (1-\epsilon)
= 1-\epsilon , \\
& U(P,f) = \sum_{i=1}^3 M_i \Delta x_i =
1 \cdot (1-\epsilon) + 1 \cdot 2\epsilon + 0 \cdot (1-\epsilon)
= 1+\epsilon .
\end{align*}
Logo,
\begin{equation*}
\overline{\int_0^2} f - 
\underline{\int_0^2} f
\leq
U(P,f) - L(P,f)
=
(1+\epsilon)
- (1-\epsilon) = 2 \epsilon .
\end{equation*}
Pela \propref{intulbound:prop}, $\underline{\int_0^2} f \leq \overline{\int_0^2} f$.
Como $\epsilon$ era arbitrário,
$\overline{\int_0^2} f = \underline{\int_0^2} f$. Portanto, $f$ é integrável
à Riemann. Por fim,
\begin{equation*}
1-\epsilon = L(P,f) \leq \int_0^2 f \leq U(P,f) =
1+\epsilon.
\end{equation*}
Logo, $\bigl\lvert \int_0^2 f - 1 \bigr\rvert \leq \epsilon$. Como
$\epsilon$ era arbitrário, concluímos que $\int_0^2 f = 1$.
\end{example}

Pode ser útil destacar parte da técnica do exemplo em uma proposição.
Observe que $U(P,f)-L(P,f)$ é exatamente a área total das partes não
sombreadas dos retângulos da \figureref{darbouxfig}.

\begin{prop}
Seja $f \colon [a,b] \to \R$ uma função limitada. Então $f$ é integrável à
Riemann se, para todo $\epsilon > 0$, existe uma partição $P$ de
$[a,b]$ tal que
\begin{equation*}
U(P,f) - L(P,f) < \epsilon .
\end{equation*}
\end{prop}

\begin{proof}
Se existe tal partição $P$ para todo $\epsilon > 0$, então
\begin{equation*}
0 \leq
\overline{\int_a^b} f - 
\underline{\int_a^b} f
\leq
U(P,f) - L(P,f) < \epsilon .
\end{equation*}
Portanto,
$\overline{\int_a^b} f = \underline{\int_a^b} f$, e $f$ é integrável.
\end{proof}

\begin{example}
Vamos mostrar que $\frac{1}{1+x}$ é integrável em $[0,b]$ para todo $b > 0$.
Mais adiante veremos que funções contínuas são integráveis, mas vamos
demonstrar diretamente como fazer isso.

Seja $\epsilon > 0$. Tome $n \in \N$ e
considere $x_i \coloneqq \nicefrac{ib}{n}$, de modo que os pontos formem a
partição $P \coloneqq \{ x_0,x_1,\ldots,x_n \}$ de $[0,b]$.
Então $\Delta x_i = \nicefrac{b}{n}$ para todo $i$.
Como $f$ é decrescente, em todo subintervalo $[x_{i-1},x_i]$,
\begin{equation*}
%mbxlatex \begin{aligned}
m_i
%mbxlatex &
=
\inf \left\{ \frac{1}{1+x} : x \in [x_{i-1},x_i] \right\} = \frac{1}{1+x_i} ,
%mbxSTARTIGNORE
\quad
%mbxENDIGNORE
%mbxlatex \\
M_i
=
%mbxlatex &
\sup \left\{ \frac{1}{1+x} : x \in [x_{i-1},x_i] \right\} =
\frac{1}{1+x_{i-1}} .
%mbxlatex \end{aligned}
\end{equation*}
Então
\begin{multline*}
U(P,f)-L(P,f)
=
\sum_{i=1}^n
\Delta x_i
(M_i-m_i)
=
\frac{b}{n}
\sum_{i=1}^n 
\left( \frac{1}{1+\nicefrac{(i-1)b}{n}} - \frac{1}{1+\nicefrac{ib}{n}} \right) 
=
\\
=
\frac{b}{n}
\left( \frac{1}{1+\nicefrac{0b}{n}} - \frac{1}{1+\nicefrac{nb}{n}} \right) 
=
\frac{b^2}{n(b+1)} .
\end{multline*}
A soma é telescópica: os termos se cancelam sucessivamente, como já vimos.
Escolhendo $n$ de modo que
$\frac{b^2}{n(b+1)} < \epsilon$, satisfazemos a proposição e, portanto, a
função é integrável.
\end{example}

\begin{remark}
Uma maneira de pensar na integral é que ela soma (integra) muitas informações
locais: soma $f(x)\,dx$ sobre todos os valores de $x$.
Leibniz escolheu o sinal de integral como um $S$ alongado para representar
uma soma.
Ao contrário das derivadas, que são \myquote{locais},
as integrais aparecem em aplicações quando se busca uma resposta
\myquote{global}: distância total percorrida, temperatura média, carga total
etc.
\end{remark}

\subsection{Mais notação}

Quando $f \colon S \to \R$ está definida em um conjunto maior $S$ e
$[a,b] \subset S$,
dizemos que $f$ é integrável à Riemann em $[a,b]$ se a restrição de $f$
a $[a,b]$ for integrável à Riemann.
Nesse caso, escrevemos $f \in \sR\bigl([a,b]\bigr)$
e usamos $\int_a^b f$ para indicar a integral de Riemann
da restrição de $f$ a $[a,b]$.

É útil definir a integral $\int_a^b f$ mesmo quando
$a \nless b$. Suponha que $b < a$ e $f \in \sR\bigl([b,a]\bigr)$.
Defina
\begin{equation*}
\int_a^b f \coloneqq - \int_b^a f .
\end{equation*}
Para qualquer função $f$, definimos
\begin{equation*}
\int_a^a f \coloneqq 0 .
\end{equation*}

Às vezes, a variável $x$ já pode ter outro significado. Quando
precisamos indicar a variável de integração, podemos simplesmente
usar outra letra. Por exemplo,
\begin{equation*}
\int_a^b f(s)\,ds \coloneqq \int_a^b f(x)\,dx .
\end{equation*}

\subsection{Exercícios}

\begin{exercise}
Defina $f \colon [0,1] \to \R$ por $f(x) \coloneqq x^3$
e seja $P \coloneqq \{ 0, 0.1, 0.4, 1 \}$. Calcule $L(P,f)$ e $U(P,f)$.
\end{exercise}

\begin{exercise}
Seja $f \colon [0,1] \to \R$ definida por $f(x) \coloneqq x$.
Mostre que $f \in \sR\bigl([0,1]\bigr)$ e
calcule $\int_0^1 f$ usando a definição de integral
(mas sinta-se à vontade para usar as proposições desta seção).%\propref{intulbound:prop}).
\end{exercise}

\begin{exercise}
Seja $f \colon [a,b] \to \R$ uma função limitada.
Suponha que exista uma sequência de partições $\{ P_k \}_{k=1}^\infty$ de $[a,b]$
tal que
\begin{equation*}
\lim_{k \to \infty} \bigl( U(P_k,f) - L(P_k,f) \bigr) = 0 .
\end{equation*}
Mostre que $f$ é integrável à Riemann e que
\begin{equation*}
\int_a^b f = 
\lim_{k \to \infty} U(P_k,f)
=
\lim_{k \to \infty} L(P_k,f) .
\end{equation*}
\end{exercise}

\begin{exercise}
Complete a demonstração da \propref{prop:refinement}.
\end{exercise}

\begin{exercise}
Suponha que $f \colon [-1,1] \to \R$ seja definida por
\begin{equation*}
f(x) \coloneqq
\begin{cases}
1 & \text{se } x > 0, \\
0 & \text{se } x \leq 0.
\end{cases}
\end{equation*}
Prove que $f \in \sR\bigl([-1,1]\bigr)$ e
calcule $\int_{-1}^1 f$ usando a definição de integral
(mas sinta-se à vontade para usar as proposições desta seção).
%(feel free to use \propref{intulbound:prop}).
\end{exercise}

\begin{exercise}
Sejam $c \in (a,b)$ e $d \in \R$.
Defina $f \colon [a,b] \to \R$ por
\begin{equation*}
f(x) \coloneqq
\begin{cases}
d & \text{se } x = c, \\
0 & \text{se } x \neq c.
\end{cases}
\end{equation*}
Prove que $f \in \sR\bigl([a,b]\bigr)$ e
calcule
$\int_a^b f$ usando a definição de integral
%(feel free to use \propref{intulbound:prop}).
(mas sinta-se à vontade para usar as proposições desta seção).
\end{exercise}

\begin{exercise} \label{exercise:taggedpartition}
Suponha que $f \colon [a,b] \to \R$ seja integrável à Riemann. Seja
$\epsilon > 0$. Mostre que existe uma partição $P = \{ x_0, x_1,
\ldots, x_n \}$
tal que, para todo
conjunto de números $\{ c_1, c_2, \ldots, c_n \}$ com
$c_k \in [x_{k-1},x_k]$ para todo $k$, temos
\begin{equation*}
\abs{\int_a^b f - \sum_{k=1}^n f(c_k) \Delta x_k} < \epsilon .
\end{equation*}
\end{exercise}

\begin{exercise}
Seja $f \colon [a,b] \to \R$ uma função integrável à Riemann.
Sejam $\alpha > 0$ e $\beta \in \R$.
Defina $g(x) \coloneqq f(\alpha x + \beta)$ no intervalo
$I = [\frac{a-\beta}{\alpha}, \frac{b-\beta}{\alpha}]$. Mostre
que $g$ é integrável à Riemann em $I$.
\end{exercise}

\begin{exercise}
Suponha que $f \colon [0,1] \to \R$ e $g \colon [0,1] \to \R$
sejam tais que $f(x) = g(x)$ para todo $x \in (0,1]$.
Suponha que $f$ seja integrável à Riemann.
Prove que $g$ é integrável à Riemann e que $\int_{0}^1 f = \int_{0}^1 g$.
\end{exercise}

\begin{exercise}
Seja $f \colon [0,1] \to \R$ uma função limitada.
Seja $P_n = \{ x_0,x_1,\ldots,x_n \}$ uma partição uniforme de $[0,1]$,
isto é, $x_i = \nicefrac{i}{n}$. A sequência $\bigl\{ L(P_n,f) \bigr\}_{n=1}^\infty$
é sempre monótona? Sim ou não: prove ou encontre um contraexemplo.
\end{exercise}

\begin{exercise}[Desafiador]
Para uma função limitada $f \colon [0,1] \to \R$, defina
$R_n \coloneqq (\nicefrac{1}{n})\sum_{i=1}^n f(\nicefrac{i}{n})$ (a
soma uniforme pela regra à direita).
\begin{enumerate}[a)]
\item
Se $f$ for integrável à Riemann, mostre que
$\int_0^1 f = \lim\limits_{n\to\infty} R_n$.
\item
Encontre uma função $f$ que não seja integrável à Riemann, mas para a qual
$\lim\limits_{n\to\infty} R_n$ exista.
\end{enumerate}
\end{exercise}

\begin{exercise}[Desafiador] \label{exercise:riemannintdarboux}
Generalize o exercício anterior.
Mostre que $f \in \sR\bigl([a,b]\bigr)$ se, e somente se, existe um $I \in \R$
tal que, para todo $\epsilon > 0$, existe
um $\delta > 0$ tal que, se $P$ é uma partição com $\Delta x_i < \delta$
para todo $i$, então
$\babs{L(P,f) - I} < \epsilon$ e
$\babs{U(P,f) - I} < \epsilon$. Se $f \in \sR\bigl([a,b]\bigr)$, então $I = \int_a^b f$.
\end{exercise}

\begin{exercise}
Usando o \exerciseref{exercise:riemannintdarboux} e a ideia da demonstração
do \exerciseref{exercise:taggedpartition}, mostre que a integral de Darboux
é equivalente à definição usual da integral de Riemann, provavelmente vista
em cálculo. Isto é, mostre que
$f \in \sR\bigl([a,b]\bigr)$ se, e somente se, existe um $I \in \R$
tal que, para todo $\epsilon > 0$, existe
um $\delta > 0$ tal que, se $P = \{ x_0,x_1,\ldots,x_n \}$
é uma partição com $\Delta x_i < \delta$
para todo $i$, então
$\abs{\sum_{i=1}^n f(c_i) \Delta x_i - I} < \epsilon$ para todo conjunto
$\{ c_1,c_2,\ldots,c_n \}$ com $c_i \in [x_{i-1},x_i]$.
Se $f \in \sR\bigl([a,b]\bigr)$, então $I = \int_a^b f$.
\end{exercise}


\begin{exercise}[Desafiador]
Construa funções $f$ e $g$ tais que
$f \colon [0,1] \to \R$ seja integrável à Riemann,
$g \colon [0,1] \to [0,1]$ seja injetiva e sobrejetiva,
mas a composição $f \circ g$ não seja integrável à Riemann.
\end{exercise}

\begin{exercise}
Suponha que $f \colon [a,b] \to \R$ seja uma função limitada e que $P$ seja
uma partição de $[a,b]$ tal que $L(P,f)=U(P,f)$. Prove que $f$ é constante.
\end{exercise}


%%%%%%%%%%%%%%%%%%%%%%%%%%%%%%%%%%%%%%%%%%%%%%%%%%%%%%%%%%%%%%%%%%%%%%%%%%%%%%

\sectionnewpage
\section{Propriedades da integral}
\label{sec:rintprop}

%mbxINTROSUBSECTION

\sectionnotes{2 aulas; pode-se omitir sem prejuízo a integrabilidade de funções
com descontinuidades}

\subsection{Aditividade}

Somar vários termos separadamente em duas partes e, em seguida, somar essas
duas partes
deveria dar o mesmo que somar tudo de uma só vez.
A propriedade correspondente para integrais é chamada de
\myindex{propriedade aditiva da integral}. Primeiro, provamos a propriedade
aditiva para as integrais inferior e superior de Darboux.

\begin{lemma} \label{lemma:darbouxadd}
Suponha que $a < b < c$ e que $f \colon [a,c] \to \R$ seja uma função limitada.
Então
\begin{equation*}
\underline{\int_a^c} f
=
\underline{\int_a^b} f
+
\underline{\int_b^c} f
\quad \text{e} \quad
\overline{\int_a^c} f
=
\overline{\int_a^b} f
+
\overline{\int_b^c} f .
\end{equation*}
\end{lemma}

\begin{proof}
Sejam $P_1 = \{ x_0,x_1,\ldots,x_k \}$
uma partição de $[a,b]$ e $P_2 = \{ x_k, x_{k+1}, \ldots, x_n \}$ uma partição
de $[b,c]$.
Então o conjunto $P \coloneqq P_1 \cup P_2 = \{ x_0, x_1, \ldots, x_n \}$ é
uma partição de $[a,c]$. Obtemos
\begin{equation*}
L(P,f) =
\sum_{i=1}^n m_i \Delta x_i
=
\sum_{i=1}^k m_i \Delta x_i
+
\sum_{i=k+1}^n m_i \Delta x_i
=
L(P_1,f) + L(P_2,f) .
\end{equation*}
Ao tomarmos o supremo do lado direito sobre todas as escolhas de $P_1$ e $P_2$,
tomamos o supremo do lado esquerdo sobre todas as partições $P$ de $[a,c]$
que contêm $b$. Se $Q$ é uma partição de $[a,c]$ e $P = Q \cup \{ b \}$,
então $P$ é um refinamento de $Q$ e, portanto, $L(Q,f) \leq L(P,f)$.
Assim, basta tomar o supremo apenas sobre as partições $P$ que contêm $b$
para obter o supremo de $L(P,f)$ sobre todas as partições $P$; veja o
\exerciseref{exercise:dominatingb}.
Por fim, lembrando o \exerciseref{exercise:supofsum},
calculamos
\begin{equation*}
\begin{split}
\underline{\int_a^c} f
& =
\sup \, \bigl\{ L(P,f) : P \text{ é uma partição de } [a,c] \bigr\}
\\
& =
\sup \, \bigl\{ L(P,f) : P \text{ é uma partição de } [a,c], b \in P \bigr\}
\\
& =
\sup \, \bigl\{ L(P_1,f) + L(P_2,f) :
P_1 \text{ é uma partição de } [a,b], P_2 \text{ é uma partição de } [b,c] \bigr\}
\\
& =
\sup \, \bigl\{ L(P_1,f) : P_1 \text{ é uma partição de } [a,b] \bigr\}
+
\sup \, \bigl\{ L(P_2,f) : P_2 \text{ é uma partição de } [b,c] \bigr\}
\\
&=
\underline{\int_a^b} f + \underline{\int_b^c} f .
\end{split}
\end{equation*}

De modo semelhante, para $P$, $P_1$ e $P_2$ como acima, obtemos
\begin{equation*}
U(P,f) =
\sum_{i=1}^n M_i \Delta x_i
=
\sum_{i=1}^k M_i \Delta x_i
+
\sum_{i=k+1}^n M_i \Delta x_i
=
U(P_1,f) + U(P_2,f) .
\end{equation*}
Queremos tomar o ínfimo do lado direito
sobre todas as escolhas de $P_1$ e $P_2$; assim, tomamos o ínfimo
sobre todas as partições $P$ de $[a,c]$ que contêm $b$. Se $Q$ é uma partição
de $[a,c]$ e $P = Q \cup \{ b \}$, então $P$ é um refinamento de $Q$
e, portanto, $U(Q,f) \geq U(P,f)$. Assim, basta tomar o ínfimo apenas sobre
as partições $P$ que contêm $b$ para obter o ínfimo de $U(P,f)$ sobre
todas as partições $P$.
Obtemos
\begin{equation*}
\overline{\int_a^c} f
=
\overline{\int_a^b} f + \overline{\int_b^c} f .  \qedhere
\end{equation*}
\end{proof}

\begin{prop}
Seja $a < b < c$. Uma função $f \colon [a,c] \to \R$ é integrável à Riemann
se, e somente se, $f$ é integrável à Riemann em $[a,b]$ e $[b,c]$. Se
$f$ for integrável à Riemann, então
\begin{equation*}
\int_a^c f
=
\int_a^b f
+
\int_b^c f .
\end{equation*}
\end{prop}

\begin{proof}
Suponha que $f \in \sR\bigl([a,c]\bigr)$. Então $f$ é limitada e
$\overline{\int_a^c} f = 
\underline{\int_a^c} f = 
\int_a^c f$. O lema fornece
\begin{equation*}
\int_a^c f
=
\underline{\int_a^c} f
 =
\underline{\int_a^b} f + \underline{\int_b^c} f
 \leq
\overline{\int_a^b} f + \overline{\int_b^c} f
 =
\overline{\int_a^c} f
 =
\int_a^c f .
\end{equation*}
Portanto, a desigualdade é uma igualdade:
\begin{equation*}
\underline{\int_a^b} f + \underline{\int_b^c} f
=
\overline{\int_a^b} f + \overline{\int_b^c} f .
\end{equation*}
Como também sabemos que
$\underline{\int_a^b} f \leq \overline{\int_a^b} f$
e
$\underline{\int_b^c} f \leq \overline{\int_b^c} f$; concluímos que
\begin{equation*}
\underline{\int_a^b} f
=
\overline{\int_a^b} f
\qquad \text{e} \qquad
\underline{\int_b^c} f
=
\overline{\int_b^c} f .
\end{equation*}
Assim, $f$ é integrável à Riemann em $[a,b]$ e $[b,c]$, e a fórmula desejada
vale.

Agora suponha que $f$ seja integrável à Riemann em $[a,b]$ e em $[b,c]$.
Novamente, $f$ é limitada, e o lema fornece
\begin{equation*}
\underline{\int_a^c} f
=
\underline{\int_a^b} f + \underline{\int_b^c} f
=
\int_a^b f + \int_b^c f
=
\overline{\int_a^b} f + \overline{\int_b^c} f
=
\overline{\int_a^c} f .
\end{equation*}
Portanto, $f$ é integrável à Riemann em $[a,c]$, e a integral é calculada
como indicado.
\end{proof}

Uma consequência imediata da aditividade é o corolário a seguir. Deixamos
os detalhes como exercício para o leitor.

\begin{cor} \label{intsubcor}
Se $f \in \sR\bigl([a,b]\bigr)$ e
$[c,d] \subset [a,b]$, então
a restrição $f|_{[c,d]}$ pertence a $\sR\bigl([c,d]\bigr)$.
\end{cor}

\subsection{Linearidade e monotonicidade}

Uma soma é uma função linear de suas parcelas. O mesmo vale para a integral.

\begin{prop}[Linearidade]
\index{linearidade da integral}\label{prop:integrallinear}
Sejam $f$ e $g$ pertencentes a $\sR\bigl([a,b]\bigr)$ e $\alpha \in \R$.
\begin{enumerate}[(i)]
\item $\alpha f$ pertence a $\sR\bigl([a,b]\bigr)$ e
\begin{equation*}
\int_a^b \alpha f(x) \,dx = \alpha \int_a^b f(x) \,dx .
\end{equation*}
\item $f+g$ pertence a $\sR\bigl([a,b]\bigr)$ e
\begin{equation*}
\int_a^b \bigl( f(x)+g(x) \bigr) \,dx = 
\int_a^b f(x) \,dx 
+
\int_a^b g(x) \,dx .
\end{equation*}
\end{enumerate}
\end{prop}

\begin{proof}
\pagebreak[2]
Vamos provar o primeiro item para $\alpha \geq 0$.
Seja $P$ uma partição de $[a,b]$, e
seja $m_i \coloneqq \inf \bigl\{ f(x) : x \in [x_{i-1},x_i] \bigr\}$, como de costume.
Como $\alpha \geq 0$, a multiplicação por $\alpha$ pode ser colocada
para fora do ínfimo:
\begin{equation*}
\inf \bigl\{ \alpha f(x) : x \in [x_{i-1},x_i] \bigr\}
=
\alpha \inf \bigl\{ f(x) : x \in [x_{i-1},x_i] \bigr\} = \alpha m_i .
\end{equation*}
Portanto,
\begin{equation*}
L(P,\alpha f) =
\sum_{i=1}^n \alpha m_i \Delta x_i = \alpha \sum_{i=1}^n m_i \Delta x_i = \alpha
L(P,f).
\end{equation*}
De modo semelhante,
\begin{equation*}
U(P,\alpha f) = \alpha U(P,f) .
\end{equation*}
Novamente, como $\alpha \geq 0$, podemos colocar a multiplicação por $\alpha$
para fora do supremo. Portanto,
\begin{equation*}
\begin{split}
\underline{\int_a^b} \alpha f(x)\,dx & =
\sup \, \bigl\{ L(P,\alpha f) : P \text{ é uma partição de } [a,b] \bigr\}
\\
& =
\sup \, \bigl\{ \alpha L(P,f) : P \text{ é uma partição de } [a,b] \bigr\}
\\
& =
\alpha \,
\sup \, \bigl\{ L(P,f) : P \text{ é uma partição de } [a,b] \bigr\}
\\
& =
\alpha
\underline{\int_a^b} f(x)\,dx .
\end{split}
\end{equation*}
De modo semelhante, mostramos que
\begin{equation*}
\overline{\int_a^b} \alpha f(x)\,dx
=
\alpha
\overline{\int_a^b} f(x)\,dx .
\end{equation*}
Agora, a conclusão segue para $\alpha \geq 0$.

Para concluir a demonstração do primeiro item (quando $\alpha < 0$), precisamos
mostrar que $-f$ é integrável à Riemann e que
$\int_a^b - f(x)\,dx =
-
\int_a^b f(x)\,dx$. Deixamos a demonstração desse fato como
\exerciseref{exercise:proofoflinpropparti}.

A demonstração do segundo item fica como
\exerciseref{exercise:proofoflinproppartii}.
O argumento não é difícil, mas tampouco é tão
trivial quanto pode parecer à primeira vista.
\end{proof}

O segundo item da proposição não vale como igualdade para as integrais de
Darboux, mas obtemos desigualdades.
A demonstração da proposição a seguir é o
\exerciseref{exercise:upperlowerlinineq}.
Para somas superiores e inferiores em uma partição fixa, ela decorre do
\exerciseref{exercise:sumofsup}, isto é, o supremo de uma soma é menor ou
igual à soma dos supremos; vale uma afirmação análoga para ínfimos.

\begin{prop} \label{prop:upperlowerlinineq}
Sejam $f \colon [a,b] \to \R$ e $g \colon [a,b] \to \R$ funções limitadas.
Então
\begin{equation*}
%\overline{\int_a^b} \bigl(f(x)+g(x)\bigr)\,dx \leq
%\overline{\int_a^b}f(x)\,dx+\overline{\int_a^b}g(x)\,dx
\overline{\int_a^b} (f+g) \leq \overline{\int_a^b}f+\overline{\int_a^b}g
,
\qquad
\text{e}
\qquad
\underline{\int_a^b} (f+g) \geq \underline{\int_a^b}f+\underline{\int_a^b}g
%\underline{\int_a^b} \bigl(f(x)+g(x)\bigr)\,dx \geq
%\underline{\int_a^b}f(x)\,dx+\underline{\int_a^b}g(x)\,dx
.
\end{equation*}
\end{prop}

Somar números menores deve produzir um resultado menor.
O mesmo vale para uma integral.

\begin{prop}[Monotonicidade]
\index{monotonicidade da integral}
Sejam $f \colon [a,b] \to \R$ e $g \colon [a,b] \to \R$
limitadas, com $f(x) \leq g(x)$
para todo $x \in [a,b]$. Então
\begin{equation*}
\underline{\int_a^b} f 
\leq
\underline{\int_a^b} g 
\qquad \text{e} \qquad
\overline{\int_a^b} f 
\leq
\overline{\int_a^b} g .
\end{equation*}
Além disso, se $f$ e $g$ pertencem a $\sR\bigl([a,b]\bigr)$, então
\begin{equation*}
\int_a^b f 
\leq
\int_a^b g .
\end{equation*}
\end{prop}

\begin{proof}
Seja $P = \{ x_0, x_1, \ldots, x_n \}$ uma partição de $[a,b]$. Defina
\begin{equation*}
m_i \coloneqq \inf \, \bigl\{ f(x) : x \in [x_{i-1},x_i] \bigr\}
\qquad \text{e} \qquad
\widetilde{m}_i \coloneqq \inf \, \bigl\{ g(x) : x \in [x_{i-1},x_i] \bigr\} .
\end{equation*}
Como $f(x) \leq g(x)$, temos $m_i \leq \widetilde{m}_i$.
Portanto,
\begin{equation*}
L(P,f)
=
\sum_{i=1}^n m_i \Delta x_i
\leq
\sum_{i=1}^n \widetilde{m}_i \Delta x_i
=
L(P,g) .
\end{equation*}
Tomamos o supremo sobre todas as partições $P$ (veja a
\propref{prop:funcsupinf}) e obtemos
\begin{equation*}
\underline{\int_a^b} f 
\leq
\underline{\int_a^b} g .
\end{equation*}
De modo semelhante, obtemos a mesma conclusão para as integrais superiores.
Por fim, se $f$ e $g$ são integráveis à Riemann, todas as integrais envolvidas
são iguais, e a conclusão segue.
\end{proof}

\subsection{Funções contínuas}

Vamos mostrar que funções contínuas são integráveis à Riemann.
Podemos até permitir algumas descontinuidades. Começamos com uma
função contínua em todo o intervalo fechado $[a,b]$.

\begin{lemma} \label{lemma:contint}
Se $f \colon [a,b] \to \R$ é contínua,
então $f \in \sR\bigl([a,b]\bigr)$.
\end{lemma}

\begin{proof}
Como $f$ é contínua em um intervalo fechado e limitado, ela é
limitada e
uniformemente contínua.
Dado $\epsilon > 0$, escolha $\delta > 0$ de modo que
$\sabs{x-y} < \delta$ implique $\babs{f(x)-f(y)} < \frac{\epsilon}{b-a}$.

Seja $P = \{ x_0, x_1, \ldots, x_n \}$
uma partição de $[a,b]$ tal que $\Delta x_i < \delta$ para todo $i = 1,2,
\ldots, n$. Por exemplo,
tome $n$ tal que $\frac{b-a}{n} < \delta$ e
defina $x_i \coloneqq \frac{i}{n}(b-a) + a$.
Então, para todos $x, y \in [x_{i-1},x_i]$, temos
$\sabs{x-y} \leq \Delta x_i < \delta$; portanto,
\begin{equation*}
f(x)-f(y) \leq \babs{f(x)-f(y)} < \frac{\epsilon}{b-a} .
\end{equation*}
Como $f$ é contínua em $[x_{i-1},x_i]$, ela atinge um máximo e um mínimo
nesse intervalo.
Seja $x$ um ponto em que $f$ atinge o máximo e seja $y$ um ponto
em que $f$ atinge o mínimo. Então $f(x) = M_i$
e $f(y) = m_i$, segundo a notação da definição da integral.
Portanto,
\begin{equation*}
M_i-m_i = f(x)-f(y) < 
\frac{\epsilon}{b-a} .
\end{equation*}
Assim,
\begin{equation*}
\begin{split}
\overline{\int_a^b} f - 
\underline{\int_a^b} f 
& \leq
U(P,f) - L(P,f)
\\
& =
\left(
\sum_{i=1}^n
M_i \Delta x_i
\right)
-
\left(
\sum_{i=1}^n
m_i \Delta x_i
\right)
\\
& =
\sum_{i=1}^n
(M_i-m_i) \Delta x_i
\\
& <
\frac{\epsilon}{b-a}
\sum_{i=1}^n
\Delta x_i
\\
& =
\frac{\epsilon}{b-a} (b-a)
= \epsilon .
\end{split}
\end{equation*}
Como $\epsilon > 0$ era arbitrário,
\begin{equation*}
\overline{\int_a^b} f = \underline{\int_a^b} f ,
\end{equation*}
e $f$ é integrável à Riemann em $[a,b]$.
\end{proof}

O segundo lema afirma que basta a função ser \myquote{integrável à Riemann no
interior do intervalo}, desde que seja limitada. Ele também nos diz como
calcular a integral.

\begin{lemma} \label{lemma:boundedimpriemann}
Seja $f \colon [a,b] \to \R$ uma função limitada, e sejam
$\{ a_n \}_{n=1}^\infty$ e $\{b_n \}_{n=1}^\infty$ sequências tais que
$a < a_n < b_n < b$ para todo $n$, com
$\lim_{n\to\infty} a_n = a$ e $\lim_{n\to\infty} b_n = b$.
Suponha que $f \in \sR\bigl([a_n,b_n]\bigr)$ para todo $n$.
Então $f \in \sR\bigl([a,b]\bigr)$ e
\begin{equation*}
\int_a^b f = 
\lim_{n \to \infty} \int_{a_n}^{b_n} f .
\end{equation*}
\end{lemma}

\begin{proof}
Seja $M > 0$ um número real tal que $\babs{f(x)} \leq M$.
Como $(b-a) \geq (b_n-a_n)$,
\begin{equation*}
-M(b-a) \leq
-M(b_n-a_n) \leq
\int_{a_n}^{b_n} f
\leq
M(b_n-a_n) \leq
M(b-a) .
\end{equation*}
Portanto, a sequência de números
$\bigl\{ \int_{a_n}^{b_n} f \bigr\}_{n=1}^\infty$ é limitada e, pelo
\hyperref[thm:bwseq]{Bolzano--Weierstrass}
possui uma subsequência convergente indexada por $n_k$. Denotemos por
$L$ o limite da subsequência
$\bigl\{ \int_{a_{n_k}}^{b_{n_k}} f \bigr\}_{k=1}^\infty$.

O \lemmaref{lemma:darbouxadd} afirma que as integrais inferior e superior
são aditivas, e a hipótese diz que $f$ é integrável em
$[a_{n_k},b_{n_k}]$. Portanto,
\begin{equation*}
\underline{\int_a^b} f
=
\underline{\int_a^{a_{n_k}}} f
+
\int_{a_{n_k}}^{b_{n_k}} f
+
\underline{\int_{b_{n_k}}^b} f
\geq
-M(a_{n_k}-a)
+
\int_{a_{n_k}}^{b_{n_k}} f
-
M(b-b_{n_k}) .
\end{equation*}
Tomamos o limite, quando $k \to \infty$, do lado direito:
\begin{equation*}
\underline{\int_a^b} f
\geq
-M\cdot 0
+
L
-
M\cdot 0
= L .
\end{equation*}

Em seguida, usamos a aditividade da integral superior:
\begin{equation*}
\overline{\int_a^b} f
=
\overline{\int_a^{a_{n_k}}} f
+
\int_{a_{n_k}}^{b_{n_k}} f
+
\overline{\int_{b_{n_k}}^b} f
\leq
M(a_{n_k}-a)
+
\int_{a_{n_k}}^{b_{n_k}} f
+
M(b-b_{n_k}) .
\end{equation*}
Tomamos o limite ao longo da mesma subsequência
$\{ \int_{a_{n_k}}^{b_{n_k}} f \}_{k=1}^\infty$ e obtemos
\begin{equation*}
\overline{\int_a^b} f
\leq
M\cdot 0
+
L
+
M\cdot 0
= L .
\end{equation*}
Assim, $\overline{\int_a^b} f = \underline{\int_a^b} f = L$;
logo, $f$ é integrável à Riemann e $\int_a^b f = L$.
Em particular, qualquer que seja a subsequência escolhida,
o valor de $L$ é o mesmo.

Para provar a afirmação final do lema, usamos a
\propref{seqconvsubseqconv:prop}. Mostramos que toda subsequência convergente
$\bigl\{ \int_{a_{n_k}}^{b_{n_k}} f \bigr\}_{k=1}^\infty$ converge a
$L = \int_a^b f$. Portanto, a sequência
$\bigl\{ \int_{a_n}^{b_n} f \bigr\}_{n=1}^\infty$ é convergente e converge a
$\int_a^b f$.
\end{proof}

Dizemos que uma função $f \colon [a,b] \to \R$ tem \emph{\myindex{um número
finito de descontinuidades}} se existe um conjunto finito
$S = \{ x_1, x_2, \ldots, x_n \} \subset [a,b]$ tal que $f$ é contínua
em todos os pontos de $[a,b] \setminus S$.

\begin{thm}
Seja $f \colon [a,b] \to \R$ uma função limitada com um número finito
de descontinuidades. Então $f \in \sR\bigl([a,b]\bigr)$.
\end{thm}

\begin{proof}
Dividimos o intervalo em um número finito de intervalos $[a_i,b_i]$
de modo que $f$ seja contínua no interior $(a_i,b_i)$. Se $f$ é contínua
em $(a_i,b_i)$, então é contínua e, portanto, integrável em $[c_i,d_i]$
sempre que $a_i < c_i < d_i < b_i$. Pelo \lemmaref{lemma:boundedimpriemann},
a restrição de $f$ a $[a_i,b_i]$ é integrável.
Pela aditividade da integral (e por \hyperref[induction:thm]{indução}),
$f$ é integrável na união desses intervalos.
\end{proof}

\subsection{Mais sobre funções integráveis}

Às vezes é conveniente (ou necessário)
alterar alguns valores de uma função e, em seguida, integrá-la. O resultado
a seguir afirma que, se alterarmos os valores em um número finito de pontos,
a integral não muda.

\begin{prop}
Seja $f \colon [a,b] \to \R$ integrável à Riemann. Seja
$g \colon [a,b] \to \R$ tal que $f(x) = g(x)$ para todo
$x \in [a,b] \setminus S$, onde $S$ é um conjunto finito.
Então $g$ é integrável à Riemann e
\begin{equation*}
\int_a^b g = \int_a^b f.
\end{equation*}
\end{prop}

\begin{proof}[Sketch of proof]
Usando a aditividade da integral, dividimos o intervalo $[a,b]$ em
intervalos menores de modo que $f(x) = g(x)$ para todo $x$, exceto talvez
nos extremos (os detalhes ficam a cargo do leitor).

Portanto, sem perda de generalidade, suponhamos que $f(x) = g(x)$ para
todo $x \in (a,b)$. A demonstração decorre do
\lemmaref{lemma:boundedimpriemann} e fica como exercício em
\exerciseref{exercise:changeendpointsintegral}.
\end{proof}

Por fim, funções monótonas (crescentes ou decrescentes) são sempre
integráveis à Riemann. A demonstração fica a cargo do leitor como parte do
\exerciseref{exercise:boundedvariationintegrable}.

\begin{prop} \label{prop:monotoneintegrable}
Seja $f \colon [a,b] \to \R$ uma função monótona. Então $f \in
\sR\bigl([a,b]\bigr)$.
\end{prop}

\subsection{Exercícios}

\begin{exercise} \label{exercise:proofoflinpropparti}
Complete a demonstração do primeiro item da
\propref{prop:integrallinear}. Seja $f \in \sR\bigl([a,b]\bigr)$.
Prove que $-f \in \sR\bigl([a,b]\bigr)$ e
\begin{equation*}
\int_a^b - f(x) \,dx = - \int_a^b f(x) \,dx .
\end{equation*}
\end{exercise}

\begin{exercise} \label{exercise:proofoflinproppartii}
Prove o segundo item da \propref{prop:integrallinear}.
Sejam $f,g \in \sR\bigl([a,b]\bigr)$.
Prove, sem usar a \propref{prop:upperlowerlinineq}, que $f+g$ pertence a
$\sR\bigl([a,b]\bigr)$ e que
\begin{equation*}
\int_a^b \bigl( f(x)+g(x) \bigr) \,dx = 
\int_a^b f(x) \,dx 
+
\int_a^b g(x) \,dx .
\end{equation*}
Sugestão: uma maneira de fazer isso é usar a \propref{prop:refinement} para encontrar
uma única partição $P$ tal que $U(P,f)-L(P,f) < \nicefrac{\epsilon}{2}$ e
$U(P,g)-L(P,g) < \nicefrac{\epsilon}{2}$.
\end{exercise}

\begin{exercise} \label{exercise:changeendpointsintegral}
Seja $f \colon [a,b] \to \R$ integrável à Riemann, e seja
$g \colon [a,b] \to \R$ tal que $f(x) = g(x)$ para todo $x \in (a,b)$.
Prove que $g$ é integrável à Riemann e que
\begin{equation*}
\int_a^b g = \int_a^b f.
\end{equation*}
\end{exercise}

\begin{exercise}
Prove o \emph{\myindex{teorema do valor médio para integrais}}:
se $f \colon [a,b] \to \R$ é contínua, então existe
$c \in [a,b]$ tal que $\int_a^b f = f(c)(b-a)$.
\end{exercise}

\begin{exercise}
Seja $f \colon [a,b] \to \R$ uma função contínua tal que $f(x) \geq 0$
para todo $x \in [a,b]$ e $\int_a^b f = 0$. Prove que $f(x) = 0$
para todo $x$.
\end{exercise}

\begin{exercise}
Seja $f \colon [a,b] \to \R$ uma função contínua
e suponha que $\int_a^b f = 0$. Prove que existe
$c \in [a,b]$ tal que $f(c) = 0$. (Compare com o exercício anterior.)
\end{exercise}

\begin{exercise}
Sejam $f \colon [a,b] \to \R$ e $g \colon [a,b] \to \R$
funções contínuas tais que $\int_a^b f = \int_a^b g$.
Mostre que existe $c \in [a,b]$ tal que $f(c) = g(c)$.
\end{exercise}

\begin{exercise}
Seja $f \in \sR\bigl([a,b]\bigr)$. Sejam $\alpha, \beta, \gamma$ números
arbitrários em $[a,b]$ (não necessariamente em alguma ordem específica). Prove
\begin{equation*}
\int_\alpha^\gamma f =
\int_\alpha^\beta f +
\int_\beta^\gamma f .
\end{equation*}
Lembre-se do significado de $\int_a^b f$ quando $b \leq a$.
\end{exercise}

\begin{exercise}
Prove o \corref{intsubcor}.
\end{exercise}

\begin{exercise} \label{exercise:easyabsint}
Suponha que $f \colon [a,b] \to \R$ seja limitada e
tenha um número finito de descontinuidades.
Mostre que, como função de $x$, a expressão $\babs{f(x)}$ é limitada e tem um número finito de
descontinuidades; portanto, é integrável à Riemann. Em seguida, mostre que
\begin{equation*}
\abs{\int_a^b f(x)\,dx} \leq \int_a^b \babs{f(x)}\,dx .
\end{equation*}
\end{exercise}

\begin{exercise}[Difícil]
Mostre que a função de
Thomae\index{função de Thomae}, ou \myindex{função pipoca}
(veja o \exampleref{popcornfunction:example}),
é integrável à Riemann. Portanto, existe uma função integrável à Riemann
que é descontínua em todos os números racionais (um conjunto denso).

Mais precisamente, defina $f \colon [0,1] \to \R$ por $f(0)=1$ e, para $x > 0$, por
\begin{equation*}
f(x) \coloneqq 
\begin{cases}
\nicefrac{1}{k} & \text{se } x=\nicefrac{m}{k}, \text{ com }
m,k \in \N \text{ e } m \text{ e } k \text{ primos entre si,} \\
0               & \text{se } x \text{ is irrational.}
\end{cases}
\end{equation*}
Mostre que $\int_0^1 f = 0$.
\end{exercise}


\begin{exnote}
Se $I \subset \R$ é um intervalo limitado, então a função
\begin{equation*}
\varphi_I(x) \coloneqq
\begin{cases}
1 & \text{se } x \in I, \\
0 & \text{caso contrário,}
\end{cases}
\end{equation*}
é chamada de \emph{\myindex{função escada elementar}}.
\end{exnote}

\begin{exercise} \label{exercise:stepfunctionintegrable}
Seja $I$ um intervalo limitado qualquer (considere todos os tipos de
intervalos: fechados, abertos e semiabertos) e sejam $a < b$. Usando apenas a
definição de integral, mostre que a função escada elementar $\varphi_I$ é
integrável em $[a,b]$ e determine a integral em termos de $a$, $b$ e dos
extremos de $I$.
\end{exercise}

\begin{exnote}
Uma função $f$ é chamada de
\emph{\myindex{função escada}} se pode ser escrita como
\begin{equation*}
f(x) = \sum_{k=1}^n \alpha_k \varphi_{I_k} (x)
\end{equation*}
para certos números reais $\alpha_1,\alpha_2, \ldots, \alpha_n$
e certos intervalos limitados $I_1,I_2,\ldots,I_n$.
\end{exnote}

\begin{exercise}
Usando o \exerciseref{exercise:stepfunctionintegrable}, mostre que uma função
escada (veja acima) é integrável em todo intervalo $[a,b]$. Além disso,
determine a integral em termos de $a$, $b$, dos extremos de $I_k$ e dos
$\alpha_k$.
\end{exercise}

\begin{exercise}
\label{exercise:boundedvariationintegrable}
Seja $f \colon [a,b] \to \R$ uma função.
\begin{enumerate}[a)]
\item
Mostre que, se $f$ é crescente, então é integrável à Riemann.
Sugestão: use uma partição uniforme, cujos subintervalos têm todos o mesmo comprimento.
\item
Use o item a) para mostrar que, se $f$ é decrescente, então é integrável à Riemann.
\item
Suponha\footnote{Dizemos que uma função $h$ assim é de
\emph{\myindex{variação limitada}}.} que $h = f-g$, onde $f$ e $g$ são
funções crescentes em $[a,b]$. Mostre que $h$ é integrável à Riemann.
\end{enumerate}
\end{exercise}

\begin{exercise}[Desafiador] \label{exercise:hardabsint}
Suponha que $f \in \sR\bigl([a,b]\bigr)$. Prove que a função que associa a $x$
o valor $\babs{f(x)}$ também é integrável à Riemann em $[a,b]$.
Em seguida, prove a mesma desigualdade do \exerciseref{exercise:easyabsint}.
\end{exercise}

\begin{exercise}  \label{exercise:upperlowerlinineq}
Suponha que $f \colon [a,b] \to \R$ e $g \colon [a,b] \to \R$
sejam limitadas.
\begin{enumerate}[a)]
\item
Mostre que
$\overline{\int_a^b} (f+g) \leq \overline{\int_a^b}f+\overline{\int_a^b}g$ e
$\underline{\int_a^b} (f+g) \geq
\underline{\int_a^b}f+\underline{\int_a^b}g$.
\item
Encontre exemplos de $f$ e $g$ para os quais
a desigualdade é estrita. Sugestão: $f$ e $g$ não devem ser integráveis à Riemann.
\end{enumerate}
\end{exercise}

\begin{exercise}
Suponha que $f \colon [a,b] \to \R$ seja contínua e que $g \colon \R \to \R$ seja
Lipschitz contínua. Defina
\begin{equation*}
h(x) \coloneqq \int_a^b g(t-x) f(t) \, dt .
\end{equation*}
Prove que $h$ é Lipschitz contínua.
\end{exercise}

\begin{exercise} \label{exercise:rieleblem}
Prove uma versão do \emph{\myindex{lema de Riemann--Lebesgue}}
(há vários resultados com esse nome):
suponha que $f \colon [a,b] \to \R$ seja contínua e defina a sequência
$\{ x_n \}_{n=1}^\infty$ por
\begin{equation*}
x_n \coloneqq \int_a^b f(t) \sin(nt) \, dt .
\end{equation*}
Prove que $\lim\limits_{n\to\infty} x_n = 0$.
\end{exercise}

%%%%%%%%%%%%%%%%%%%%%%%%%%%%%%%%%%%%%%%%%%%%%%%%%%%%%%%%%%%%%%%%%%%%%%%%%%%%%%

\sectionnewpage
\section{Fundamental theorem of calculus}
\label{sec:ftc}

%mbxINTROSUBSECTION

\sectionnotes{1.5 lectures}

In this section we discuss and prove the
\emph{\myindex{fundamental theorem of calculus}}.
The entirety of integral calculus is built upon this theorem,
ergo the name.
The theorem relates the seemingly unrelated concepts of integral and
derivative.  It tells us how to compute the antiderivative of a function
using the integral and vice versa.

\subsection{First form of the theorem}

\begin{thm} \label{thm:FTCv1}
Let $F \colon [a,b] \to \R$ be a continuous function, differentiable
on $(a,b)$.  Let $f \in \sR\bigl([a,b]\bigr)$ be such that $f(x) = F'(x)$ for $x \in
(a,b)$.  Then
\begin{equation*}
\int_a^b f = F(b)-F(a) .
\end{equation*}
\end{thm}

It is not hard to generalize the theorem to allow a finite number of points
in $[a,b]$ where $F$ is not differentiable, as long as it is continuous.
This generalization is left as an exercise.

\begin{proof}
Let $P = \{ x_0, x_1, \ldots, x_n \}$ be a partition of $[a,b]$.
For each interval $[x_{i-1},x_i]$, use the
\hyperref[thm:mvt]{mean value theorem} to find a
$c_i \in (x_{i-1},x_i)$ such that
\begin{equation*}
f(c_i) \Delta x_i = F'(c_i) (x_i - x_{i-1}) = F(x_i) - F(x_{i-1}) .
\end{equation*}
See \figureref{fig:fundthmfig}, and
note that the area of the
$i$th rectangle is
$F(x_{i})-F(x_{i-1})$,
and the total area of all three rectangles pictured is
$F(x_{i+1})-F(x_{i-2})$.
The idea is that taking smaller and smaller subintervals,
the total area of all these rectangles converges to the integral of $f$.
\begin{myfigureht}
\myincludegraphics{fundthmfig}{%
A diagram of a graph of a function in dark bold
marked as y equals f of x equals F prime of x
and three subintervals of the x coordinates.
The middle subinterval is labeled as going from
x sub quantity i minus 1 to x sub i and is of length Delta x sub i.
A point c sub i is marked inside this subinterval and a
dashed line goes up vertically until f of c sub i where it hits
the graph of f.  A shaded rectangle of this height and the subinterval
as the base is drawn and labeled with 'area equals f of c sub i times Delta
x sub i equals F of x sub i minus F of x sub quantity i minus 1'.
The other two subintervals are similar except on the left
we replace i with i minus 1 and on the right we replace i with
i plus 1.}
\caption{Mean value theorem on subintervals of a partition
approximating the area under the curve.\label{fig:fundthmfig}}
\end{myfigureht}

Using the notation from the definition of the integral,
$m_i \leq f(c_i) \leq M_i$, and multiplying by $\Delta x_i$ gets
\begin{equation*}
m_i \Delta x_i \leq F(x_i) - F(x_{i-1}) \leq M_i \Delta x_i .
\end{equation*}
We sum over $i = 1,2, \ldots, n$ to get
\begin{equation*}
\sum_{i=1}^n m_i \Delta x_i
\leq \sum_{i=1}^n \bigl(F(x_i) - F(x_{i-1}) \bigr)
\leq \sum_{i=1}^n M_i \Delta x_i .
\end{equation*}
In the middle sum, all the terms except the first and last cancel 
and we end up with $F(x_n)-F(x_0) = F(b)-F(a)$.  The sums on the left
and on the right are the lower and the upper sums, respectively.  So
\begin{equation*}
L(P,f) \leq F(b)-F(a) \leq U(P,f) .
\end{equation*}
We take the supremum of $L(P,f)$ over all partitions $P$ and the left inequality
yields 
\begin{equation*}
\underline{\int_a^b} f \leq F(b)-F(a) .
\end{equation*}
Similarly, taking
the infimum of $U(P,f)$ over all partitions $P$ yields
\begin{equation*}
F(b)-F(a) \leq \overline{\int_a^b} f .
\end{equation*}
As $f$ is Riemann integrable, we have
\begin{equation*}
\int_a^b f =
\underline{\int_a^b} f \leq F(b)-F(a) \leq \overline{\int_a^b} f
= \int_a^b f .
\end{equation*}
The inequalities must be equalities and we are done.
\end{proof}

The theorem is used to compute integrals.  Suppose we know that
the function $f(x)$ is a derivative of some other function $F(x)$,
then we can find an explicit expression for $\int_a^b f$. 

\begin{example}
To compute
\begin{equation*}
\int_0^1 x^2 \,dx ,
\end{equation*}
we notice $x^2$ is the derivative of $\frac{x^3}{3}$.
The fundamental theorem says
\begin{equation*}
\int_0^1 x^2 \,dx =
\frac{1^3}{3}
-
\frac{0^3}{3}
= \frac{1}{3}.
\end{equation*}
\end{example}

\subsection{Second form of the theorem}

The second form of the fundamental theorem gives us a way to solve
the differential equation $F'(x) = f(x)$, where $f$ is a known
function and we are trying to find an $F$ that satisfies the equation.

\begin{thm} \label{thm:FTCv2}
Let $f \colon [a,b] \to \R$ be a Riemann integrable function.  Define
\begin{equation*}
F(x) \coloneqq \int_a^x f .
\end{equation*}
First, $F$ is continuous on $[a,b]$.  Second,
if $f$ is continuous at $c \in [a,b]$, then $F$ is differentiable at $c$
and $F'(c) = f(c)$.
\end{thm}

\begin{proof}
As $f$ is bounded, there is an $M > 0$
such that $\babs{f(x)} \leq M$ for all $x \in [a,b]$.  Suppose $x,y \in [a,b]$
with $x > y$.  Then
\begin{equation*}
\babs{F(x)-F(y)} =
\abs{\int_a^x f - \int_a^y f}
=
\abs{\int_y^x f}
\leq
M\sabs{x-y} .
\end{equation*}
By symmetry, the same also holds if $x < y$.
So $F$ is Lipschitz continuous and hence continuous.

Now suppose $f$ is continuous at $c$.
Let $\epsilon > 0$ be given.  Let $\delta > 0$ be such that
for $x \in [a,b]$,
$\sabs{x-c} < \delta$ implies $\babs{f(x)-f(c)} < \epsilon$.
In particular,
for such $x$, we have
\begin{equation*}
f(c)-\epsilon < f(x) < f(c) + \epsilon.
\end{equation*}
Thus if $x > c$, then
\begin{equation*}
\bigl(f(c)-\epsilon\bigr) (x-c) \leq \int_c^x f \leq
\bigl(f(c) + \epsilon\bigr)(x-c).
\end{equation*}
When $c > x$, then the inequalities are reversed.  Therefore,
assuming $x \neq c$, we get
\begin{equation*}
f(c)-\epsilon
\leq
\frac{\int_c^{x} f}{x-c}
\leq
f(c)+\epsilon .
\end{equation*}
As 
\begin{equation*}
\frac{F(x)-F(c)}{x-c}
=
\frac{\int_a^{x} f - \int_a^{c} f}{x-c}
=
\frac{\int_c^{x} f}{x-c} ,
\end{equation*}
we have 
\begin{equation*}
\abs{\frac{F(x)-F(c)}{x-c} - f(c)} \leq \epsilon .
\end{equation*}
The result follows.  It is left to the reader to see why is it OK that we
just have a non-strict inequality.
\end{proof}

Of course, if $f$ is continuous on $[a,b]$, then it is automatically Riemann
integrable, $F$ is differentiable on all of $[a,b]$ and $F'(x) = f(x)$ for
all $x \in [a,b]$.

\begin{remark} \label{remark:fundthmbase}
The second form of the fundamental theorem of calculus still holds if
we let $d \in [a,b]$ and define
\begin{equation*}
F(x) \coloneqq \int_d^x f .
\end{equation*}
That is, we can use any point of $[a,b]$ as our base point.  The proof is
left as an exercise.
\end{remark}

Let us look at what a simple discontinuity can do.  Take $f(x) \coloneqq -1$ if $x
< 0$, and $f(x) \coloneqq 1$ if $x \geq 0$.  Let $F(x) \coloneqq \int_0^x f$.  It is not
difficult to see that $F(x) = \sabs{x}$.  Notice that $f$ is discontinuous at
$0$ and $F$ is not differentiable at $0$.  However, the converse in the
theorem does not
hold.
Let $g(x) \coloneqq 0$ if $x \neq 0$, and $g(0) \coloneqq 1$.
Letting $G(x) \coloneqq \int_0^x g$,
we find that $G(x) = 0$ for all $x$.  So $g$ is discontinuous
at $0$, but $G'(0)$ exists and is equal to 0.

A common misunderstanding of the integral for calculus students is to
think of integrals whose solution cannot be given in closed-form as somehow
deficient.  This is not the case.  Most integrals we write down are not
computable in closed-form.  Even some integrals that we consider
in closed-form are not really such.  We define
the natural logarithm as the antiderivative of $\nicefrac{1}{x}$
such that $\ln 1 = 0$:
\begin{equation*}
\ln x \coloneqq \int_1^x \frac{1}{s}\,ds .
\end{equation*}
How does a computer find the value of $\ln x$?
One way to do it is to numerically approximate this integral.
Morally,
we did not really \myquote{simplify} $\int_1^x \frac{1}{s}\,ds$ by
writing down $\ln x$.  We simply gave the integral a name.
If we require numerical answers,
it is possible we end up doing
the calculation by approximating an integral anyway.
In the next section, we even define the exponential using
the logarithm, which we define in terms of the integral.

Another common function defined by an integral that cannot
be evaluated symbolically in terms of elementary functions
is the $\operatorname{erf}$ function, defined as
\begin{equation*}
\operatorname{erf}(x) \coloneqq \frac{2}{\sqrt{\pi}} \int_0^x e^{-s^2} \,ds .
\end{equation*}
This function comes up often in applied mathematics.  It is simply 
the antiderivative of $\left(\nicefrac{2}{\sqrt{\pi}}\right) e^{-x^2}$
that is zero at zero.
The second form of the fundamental theorem tells us that we can write the function
as an integral.  If we wish to compute any particular value, we 
numerically approximate the integral.

\subsection{Change of variables}

A theorem often used in calculus to solve integrals is the change of
variables theorem, you may have called it
\emph{$u$-substitution}\index{u-substitution@$u$-substitution}.
To make it easy to connect to your calculus class, we will also
call the new variable $u$, that is, we make the substitution $u=g(x)$.
Recall a function is continuously differentiable if
it is differentiable and the derivative is continuous.

\begin{thm}[Change of variables]
\index{change of variables theorem}
Let $g \colon [a,b] \to \R$ be a continuously differentiable function,
let $f \colon [c,d] \to \R$ be continuous, and suppose
$g\bigl([a,b]\bigr) \subset [c,d]$.  Then
\begin{equation*}
\int_a^b f\bigl(g(x)\bigr)\, g'(x)\, dx =
\int_{g(a)}^{g(b)} f(u)\, du .
\end{equation*}
\end{thm}

\begin{proof}
As $g$, $g'$, and $f$ are continuous, $f\bigl(g(x)\bigr)\,g'(x)$
is a continuous function of $[a,b]$, therefore it is Riemann integrable.
Similarly, $f$ is integrable on every subinterval of $[c,d]$.

Define $F \colon [c,d] \to \R$ by
\begin{equation*}
F(y) \coloneqq \int_{g(a)}^{y} f(u)\,du .
\end{equation*}
By the second form of the fundamental
theorem of calculus
(see \remarkref{remark:fundthmbase} and \exerciseref{secondftc:exercise}),
$F$ is a differentiable function and $F'(y) = f(y)$.  Apply the chain
rule,
\begin{equation*}
\bigl( F \circ g \bigr)' (x) =
F'\bigl(g(x)\bigr) g'(x)
=
f\bigl(g(x)\bigr) g'(x) .
\end{equation*}
Note that $F\bigl(g(a)\bigr) = 0$ and
use the first form of the fundamental theorem
to obtain
\begin{multline*}
%mbxSTARTIGNORE
\qquad %to center things more
%mbxENDIGNORE
\int_{g(a)}^{g(b)} f(u)\,du = F\bigl(g(b)\bigr) = F\bigl(g(b)\bigr)-F\bigl(g(a)\bigr)
\\
=
\int_a^b 
\bigl( F \circ g \bigr)' (x) \,dx
=
\int_a^b 
f\bigl(g(x)\bigr) g'(x)
\,dx .
%mbxSTARTIGNORE
\qquad %to center things more
%mbxENDIGNORE
\qedhere
\end{multline*}
\end{proof}

The change of variables theorem is often used to solve integrals by changing them
to integrals that we know or that we can solve using the fundamental theorem of
calculus.

\begin{example}
The derivative of $\sin(x)$ is $\cos(x)$.
Using $g(x) \coloneqq x^2$, we solve
\begin{equation*}
\int_0^{\sqrt{\pi}} x \cos(x^2) \, dx = \int_0^\pi \frac{\cos(u)}{2} \, du
=
\frac{1}{2}
\int_0^\pi \cos(u) \, du
=
\frac{
\sin(\pi) - \sin(0)
}{2}
=
0 .
\end{equation*}
\end{example}

However, beware that we must satisfy the hypotheses of the theorem.  The
following example demonstrates why we should not just 
move symbols around mindlessly.
We must be careful that those symbols really make sense.

\begin{example}
Consider
\begin{equation*}
\int_{-1}^{1} \frac{\ln \sabs{x}}{x} \,dx .
\end{equation*}
It may be tempting to take $g(x) \coloneqq \ln \sabs{x}$.  Compute $g'(x) =
\nicefrac{1}{x}$ and try to write
\begin{equation*}
\int_{g(-1)}^{g(1)} u \,du = 
\int_{0}^{0} u \,du = 0. 
\end{equation*}
This \myquote{solution} is incorrect, and it does not say
that we can solve the given integral.  The first problem is that
$\frac{\ln \sabs{x}}{x}$ is not continuous on $[-1,1]$.
It is not defined at 0, and cannot be made continuous by defining a value at
0.
Second, $\frac{\ln \sabs{x}}{x}$ is not even Riemann integrable on $[-1,1]$
(it is unbounded).
The integral we wrote down simply does not make sense.
Finally, $g$ is not continuous 
on $[-1,1]$, let alone continuously differentiable.
\end{example}

\subsection{Exercícios}

\begin{exercise}
Compute
$\displaystyle
\frac{d}{dx} \biggl( \int_{-x}^x e^{s^2}\,ds \biggr)$.
\end{exercise}

\begin{exercise}
Compute
$\displaystyle
\frac{d}{dx} \biggl( \int_{0}^{x^2} \sin(s^2)\,ds \biggr)$.
\end{exercise}

\begin{exercise}
Suppose $F \colon [a,b] \to \R$ is continuous and differentiable
on $[a,b] \setminus S$, where $S$ is a finite set.  Suppose there
exists an $f \in \sR\bigl([a,b]\bigr)$ such that $f(x) = F'(x)$ for $x \in [a,b]
\setminus S$.  Show that
$\int_a^b f = F(b)-F(a)$.
\end{exercise}

\begin{exercise} \label{secondftc:exercise}
Let $f \colon [a,b] \to \R$ be a continuous function.  Let $c \in [a,b]$
be arbitrary.  Define
\begin{equation*}
F(x) \coloneqq \int_c^x f .
\end{equation*}
Prove that $F$ is differentiable and that $F'(x) = f(x)$ for all $x \in
[a,b]$.
\end{exercise}

\begin{exercise}
Prove \emph{\myindex{integration by parts}}.  That is, suppose $F$ and
$G$ are continuously differentiable functions on $[a,b]$.  Then prove
\begin{equation*}
\int_a^b F(x)G'(x)\,dx
=
F(b)G(b)-F(a)G(a)
-
\int_a^b F'(x)G(x)\,dx .
\end{equation*}
\end{exercise}

\begin{exercise}
Suppose $F$ and $G$ are
continuously\footnote{
Compare this hypothesis to \exerciseref{exercise:samediffconst}.}
differentiable
functions defined on $[a,b]$
such that $F'(x) = G'(x)$ for all $x \in [a,b]$.
Using the fundamental theorem of calculus,
show that $F$ and $G$ differ by a constant.  That is, show that
there exists a $C \in \R$ such that
$F(x)-G(x) = C$.
\end{exercise}

\begin{exnote}
The next exercise shows how we can use the integral to \myquote{smooth out} a
non-differentiable function.
\end{exnote}

\begin{exercise} \label{exercise:smoothingout}
Let $f \colon [a,b] \to \R$ be a continuous function.  Let $\epsilon > 0$
be a constant such that $a+\epsilon < b-\epsilon$.  For $x \in [a+\epsilon,b-\epsilon]$, define
\begin{equation*}
g(x) \coloneqq \frac{1}{2\epsilon} \int_{x-\epsilon}^{x+\epsilon} f .
\end{equation*}
\begin{enumerate}[a)]
\item
Show that $g$ is differentiable and find the derivative.
\item
Let $f$ be differentiable and fix $x \in (a,b)$ (let $\epsilon$
be small enough).  What happens to $g'(x)$ as $\epsilon$ gets smaller?
\item
Find $g$ for $f(x) \coloneqq \sabs{x}$, $\epsilon = 1$ (you can assume 
$[a,b]$ is large enough).
\end{enumerate}
\end{exercise}

\begin{exercise}
Suppose $f \colon [a,b] \to \R$ is continuous and
$\int_a^x f = \int_x^b f$ for all $x \in [a,b]$.  Show that $f(x) = 0$
for all $x \in [a,b]$.
\end{exercise}

\begin{exercise}
Suppose $f \colon [a,b] \to \R$ is continuous and
$\int_a^x f = 0$ for all rational $x$ in $[a,b]$.  Show that $f(x) = 0$
for all $x \in [a,b]$.
\end{exercise}

\begin{samepage}
\begin{exercise}
A function $f$ is an \emph{\myindex{odd function}} if $f(x) = -f(-x)$,
and $f$ is an \emph{\myindex{even function}} if $f(x) = f(-x)$.  Let $a >
0$.  Assume $f$ is continuous.  Prove:
\begin{enumerate}[a)]
\item
If $f$ is odd, then $\int_{-a}^a f = 0$.
\item
If $f$ is even, then $\int_{-a}^a f = 2 \int_0^a f$.
\end{enumerate}
\end{exercise}
\end{samepage}

\begin{exercise}
\leavevmode
\begin{enumerate}[a)]
\item
Show that $f(x) \coloneqq \sin(\nicefrac{1}{x})$
is integrable on every interval (you can define $f(0)$ to be anything).
\item
Compute $\int_{-1}^1 \sin(\nicefrac{1}{x})\,dx$ (mind the
discontinuity).
\end{enumerate}
\end{exercise}

\begin{exercise}[uses \sectionref{sec:monotonefunc}]
\leavevmode
\begin{enumerate}[a)]
\item
Suppose $f \colon [a,b] \to \R$ is increasing.
By \propref{prop:monotoneintegrable},
%\exerciseref{exercise:boundedvariationintegrable},
$f$ is Riemann integrable.  Show that if $f$ has a discontinuity at $c \in
(a,b)$, then $F(x) \coloneqq \int_a^x f$ is not differentiable at $c$.
\item
In \exerciseref{exercise:increasingfuncdiscatQ}, you constructed an increasing
function $f \colon [0,1] \to \R$ that is discontinuous at every
$x \in [0,1] \cap \Q$.  Use this $f$ to construct a function $F(x)$ that is
continuous on $[0,1]$, but not differentiable at all $x \in [0,1] \cap \Q$.
\end{enumerate}
\end{exercise}

\begin{exercise}
For any $\ell \in \N$,
show that the following limit exists and find what it is:
\begin{equation*}
\lim_{n\to\infty} \sum_{k=1}^n \frac{k^\ell}{n^{\ell+1}}
\end{equation*}
\end{exercise}

%%%%%%%%%%%%%%%%%%%%%%%%%%%%%%%%%%%%%%%%%%%%%%%%%%%%%%%%%%%%%%%%%%%%%%%%%%%%%%

\sectionnewpage
\section{The logarithm and the exponential}
\label{sec:logandexp}

%mbxINTROSUBSECTION

\sectionnotes{1 lecture (optional, requires the optional sections 
\sectionref{sec:limitatinf},
\sectionref{sec:monotonefunc},
\sectionref{sec:ift})}

We now have the tools required to properly define the exponential
and the
logarithm that you know from calculus so well.  We start with
exponentiation.
If $n$ is a positive integer, we define
\begin{equation*}
x^n \coloneqq \underbrace{x \cdot x \cdots x}_{n \text{ times}} .
\end{equation*}
It makes sense to define $x^0 \coloneqq 1$.
For negative integers, let $x^{-n} \coloneqq \nicefrac{1}{x^n}$ as long as
$x \neq 0$.
Next suppose $x > 0$.
Define
$x^{1/n}$ as
the unique positive $n$th root.  Finally, for a rational
number $\nicefrac{n}{m}$ (in lowest terms), define
\begin{equation*}
x^{n/m} \coloneqq {\bigl(x^{1/m}\bigr)}^n .
\end{equation*}
It is not difficult to show 
we get the same number no matter what
representation of $\nicefrac{n}{m}$ we use, so we do not need to use
lowest terms.

However, what do we mean by $\sqrt{2}^{\sqrt{2}}$?  Or
$x^y$ in general?  In particular, what is $e^x$ for all~$x$?
And how do we solve $y=e^x$ for~$x$?
This section answers these questions and more.

\subsection{The logarithm}
\index{logarithm}

It is convenient to define the logarithm first.
Let us show that
a unique function with the right properties exists, and only then will
we call it \emph{the} logarithm.

\begin{prop}
There exists a unique function $L \colon (0,\infty) \to \R$ such that
\begin{enumerate}[(i)]
\item \label{it:log:i}
$L(1) = 0$.
\item \label{it:log:ii}
$L$ is differentiable and $L'(x) = \nicefrac{1}{x}$.
\item \label{it:log:iii}
$L$ is strictly increasing, bijective, and
\begin{equation*}
\lim_{x\to 0} L(x) = -\infty , \qquad \text{e} \qquad
\lim_{x\to \infty} L(x) = \infty .
\end{equation*}
\item \label{it:log:iv}
$L(xy) = L(x)+L(y)$ for all $x,y \in (0,\infty)$.
\item \label{it:log:v}
If $q$ is a rational number and $x > 0$, then
$L(x^q) = q L(x)$.
\end{enumerate}
\end{prop}

\begin{proof}
To prove existence, we define a candidate and show it satisfies
all the properties.  Let
\begin{equation*}
L(x) \coloneqq \int_1^x \frac{1}{t}\,dt .
\end{equation*}

Obviously, \ref{it:log:i} holds.  Property \ref{it:log:ii} holds
via the second form of the fundamental theorem of calculus
(\thmref{thm:FTCv2}).

To prove property \ref{it:log:iv},
we change variables $u=yt$ to obtain
\begin{equation*}
L(x) =
\int_1^{x} \frac{1}{t}\,dt
=
\int_y^{xy} \frac{1}{u}\,du
=
\int_1^{xy} \frac{1}{u}\,du
-
\int_1^{y} \frac{1}{u}\,du
=
L(xy)-L(y) .
\end{equation*}

Let us prove \ref{it:log:iii}.
Property \ref{it:log:ii} together with the fact that $L'(x) = \nicefrac{1}{x} > 0$ 
for $x > 0$, implies that $L$
is strictly increasing and hence one-to-one.
Let us show $L$ is onto.  
As $\nicefrac{1}{t} \geq \nicefrac{1}{2}$ when $t \in [1,2]$,
\begin{equation*}
L(2) = \int_1^2 \frac{1}{t} \,dt \geq \nicefrac{1}{2} .
\end{equation*}
By induction, \ref{it:log:iv} implies that for $n \in \N$,
\begin{equation*}
L(2^n) = L(2) + L(2) + \cdots + L(2) = n L(2) .
\end{equation*}
Given $y > 0$, 
by the \hyperref[thm:arch:i]{Archimedean property} of the real numbers
(notice $L(2) > 0$), there is an $n \in \N$ such that
$L(2^n) > y$.  The
\hyperref[IVT:thm]{intermediate value theorem}
gives an $x_1 \in (1,2^n)$ such that $L(x_1) = y$.  Thus
$(0,\infty)$ is in the image of $L$.
As $L$ is increasing, $L(x) > y$ for all $x > 2^n$, and so
\begin{equation*}
\lim_{x\to\infty} L(x) = \infty .
\end{equation*}
Next
$0 = L(\nicefrac{x}{x}) = L(x) + L(\nicefrac{1}{x})$, and
so $L(x) = - L(\nicefrac{1}{x})$.  Using $x=2^{-n}$, we obtain
as above that $L$ achieves all negative numbers.  And
\begin{equation*}
\lim_{x \to 0} L(x) = 
\lim_{x \to 0} -L(\nicefrac{1}{x})
=
\lim_{x \to \infty} -L(x)
=  - \infty .
\end{equation*}
In the limits, note that only $x > 0$ are in the domain of $L$.

Let us prove \ref{it:log:v}.
Fix $x > 0$.
As above, \ref{it:log:iv} implies
$L(x^n) = n L(x)$
for all $n \in \N$.
We already found that
$L(x) = - L(\nicefrac{1}{x})$,
so $L(x^{-n}) = - L(x^n) = -n L(x)$.  Then for $m \in \N$
\begin{equation*}
L(x) = L\Bigl({(x^{1/m})}^m\Bigr) = m L\bigl(x^{1/m}\bigr) .
\end{equation*}
Putting everything together for $n \in \Z$ and $m \in \N$, we have
$L(x^{n/m}) = n L(x^{1/m}) = (\nicefrac{n}{m}) L(x)$.

Uniqueness follows using properties \ref{it:log:i} and
\ref{it:log:ii}.  Via the first form of the
fundamental theorem of calculus (\thmref{thm:FTCv1}),
\begin{equation*}
L(x) = \int_1^x \frac{1}{t}\,dt
\end{equation*}
is the unique function such that $L(1) = 0$ and $L'(x) = \nicefrac{1}{x}$.
\end{proof}

Having proved
that there is a unique function with these properties,
we simply define the \emph{\myindex{logarithm}}, sometimes called the
\emph{\myindex{natural logarithm}}:
\glsadd{not:ln}
\begin{equation*}
\ln(x) \coloneqq L(x) .
\end{equation*}
See \figureref{fig:log}.
Mathematicians usually write $\log(x)$ instead of $\ln(x)$, which is
more familiar to calculus students.  For all practical purposes, there is only one logarithm:
the natural logarithm.  See \exerciseref{exercise:otherlogbases}.
\begin{myfigureht}
\myincludegraphics{logfig}{%
A graph y equals ln of x in a bold line starting steeply
on the left in the 
negative values and crossing the x-axis at 1, and getting less and less
steep as we go further right.
A plot of y equals 1 over x is in dashed line, it is a line
that is always positive starting high on the left and getting
less and less steep as we go to the right.
The area between the dashed line and the x-axis
for x between 1 and 4 is shaded and the height of the graph
of ln at 4 is marked as 'shaded area equals ln of 4'.}
\caption{Plot of $\ln(x)$ together with $\nicefrac{1}{x}$, showing the value
$\ln(4)$.\label{fig:log}}
\end{myfigureht}

\subsection{The exponential}
\index{exponential}

Just as with the logarithm we define the exponential via a list of
properties.

\begin{prop}
There exists a unique function $E \colon \R \to (0,\infty)$ such that
\begin{enumerate}[(i)]
\item \label{it:exp:i}
$E(0) = 1$.
\item \label{it:exp:ii}
$E$ is differentiable and $E'(x) = E(x)$.
\item \label{it:exp:iii}
$E$ is strictly increasing, bijective, and
\begin{equation*}
\lim_{x\to -\infty} E(x) = 0 \qquad \text{e} \qquad
\lim_{x\to \infty} E(x) = \infty .
\end{equation*}
\item \label{it:exp:iv}
$E(x+y) = E(x)E(y)$ for all $x,y \in \R$.
\item \label{it:exp:v}
If $q \in \Q$, then
$E(qx) = {E(x)}^q$.
\end{enumerate}
\end{prop}

\begin{proof}
Again, we prove existence of such a function by defining a candidate
and proving that it satisfies all the properties.
The $L = \ln$ defined above is invertible.  Let $E$ be the
inverse function of $L$.  Property \ref{it:exp:i} is immediate.

Property \ref{it:exp:ii} follows
via the inverse function theorem, in particular
via \lemmaref{lemma:ift}:  $L$~satisfies
all the hypotheses of the lemma, and hence
\begin{equation*}
E'(x) = \frac{1}{L'\bigl(E(x)\bigr)} = E(x) .
\end{equation*}

Let us look at property \ref{it:exp:iii}.
The function $E$ is strictly increasing since 
$E'(x) = E(x) > 0$.  As $E$ is the inverse of $L$, it must also
be bijective.  
To find the limits, we use that 
$E$ is strictly increasing and onto $(0,\infty)$.
For every $M > 0$, there is an $x_0$ such that
$E(x_0) = M$ and $E(x) \geq M$ for all $x \geq x_0$.
Similarly, for every $\epsilon > 0$, there is
an $x_0$ such that $E(x_0) = \epsilon$ and
$E(x) < \epsilon$ for all $x < x_0$.
Therefore,
\begin{equation*}
\lim_{x\to -\infty} E(x) = 0 \qquad \text{e} \qquad
\lim_{x\to \infty} E(x) = \infty .
\end{equation*}

To prove property \ref{it:exp:iv}, we use the corresponding
property for the logarithm.
Take $x, y \in \R$.
As $L$ is bijective, find $a$ and $b$ such that $x = L(a)$ and $y = L(b)$.  Then
\begin{equation*}
E(x+y) =
E\bigl(L(a)+L(b)\bigr) = 
E\bigl(L(ab)\bigr) = ab = E(x)E(y)  .
\end{equation*}

Property \ref{it:exp:v} also follows from the corresponding property of $L$.
Given $x \in \R$, let $a$ be such that $x = L(a)$ and
\begin{equation*}
E(qx) = E\bigl(qL(a)\bigr)
=
E\bigl(L(a^q)\bigr) = a^q = {E(x)}^q .
\end{equation*}

Uniqueness follows from
\ref{it:exp:i} and
\ref{it:exp:ii}.
Let $E$ and $F$
be two functions satisfying
\ref{it:exp:i} and \ref{it:exp:ii}.  
\begin{equation*}
\frac{d}{dx} \Bigl( F(x)E(-x) \Bigr)
=
F'(x)E(-x) - E'(-x)F(x)
=
F(x)E(-x) - E(-x)F(x) = 0 .
\end{equation*}
Therefore, by \propref{prop:derzeroconst},
$F(x)E(-x) = F(0)E(-0) = 1$ for all $x \in \R$.
Doing the computation with $F = E$,
we obtain $E(x)E(-x) = 1$.
Then
\begin{equation*}
0 = 1-1 = F(x)E(-x) - E(x)E(-x) = \bigl(F(x)-E(x)\bigr) E(-x) .
\end{equation*}
%Since $E(x)E(-x) = 1$,
Finally, $E(-x) \neq 0$\footnote{%
$E$ is a function into $(0,\infty)$ after all.
However, $E(-x) \neq 0$ also follows
from $E(x)E(-x) = 1$.  Therefore, we can prove uniqueness of $E$ 
given \ref{it:exp:i} and \ref{it:exp:ii}, even for functions $E \colon \R
\to \R$.}
for all $x \in \R$.
So
$F(x)-E(x) = 0$ for all $x$, and we are done.
\end{proof}

Having proved $E$ is unique, we define the
\emph{\myindex{exponential}} function (see \figureref{fig:exp}) as
\glsadd{not:exp}
\begin{equation*}
\exp(x) \coloneqq E(x) .
\end{equation*}
\begin{myfigureht}
\myincludegraphics{expfig}{%
A slope field graphed on the xy-plane for x between roughly minus 1 to 1
and y between 0 and 4.  The slope field is a grid of short line segments of
a given slope.  In this case the slopes are the same in every row.
They are all upwards (positive) and
get progressively steeper as we move up.  The graph of the exponential
function is given.  It is an increasing function coming from the left
and it follows the slopes given:
Starting above the x-axis, it passes through the point (0,1), and then gets
steeper and steeper as it exits the picture on the right.}
\caption{Plot of $e^x$, together with a slope field giving slope $y$ at
every point $(x,y)$.
The equation $\frac{d}{dx} e^x = e^x$ means that $y=e^x$ follows these slopes.\label{fig:exp}}
\end{myfigureht}

If $y \in \Q$ and $x > 0$, then
\begin{equation*}
x^y = \exp\bigl(\ln(x^y)\bigr) = \exp\bigl(y\ln(x)\bigr) .
\end{equation*}
We can now make sense of exponentiation $x^y$ for arbitrary $y \in \R$;
if $x > 0$ and $y$ is irrational, define
\glsadd{not:pow}
\begin{equation*}
x^y \coloneqq \exp\bigl(y\ln(x)\bigr) .
\end{equation*}
As $\exp$ is continuous, $x^y$ is a continuous function of $y$.
Therefore, we would
obtain the same result had we taken a sequence of rational numbers
$\{ y_n \}_{n=1}^\infty$
approaching $y$ and defined $x^y = \lim_{n\to\infty} x^{y_n}$.

Define the number $e$,
called \emph{\myindex{Euler's number}} or
the \emph{\myindex{base of the natural logarithm}}, as
\glsadd{not:e}
\begin{equation*}
e \coloneqq \exp(1) .
\end{equation*}
Let us justify the notation $e^x$ for $\exp(x)$:
\begin{equation*}
e^x = \exp\bigl(x \ln(e) \bigr) = \exp(x) .
\end{equation*}

The properties of the logarithm and the exponential extend to
irrational powers.  The proof is immediate.

\begin{prop}
Let $x, y \in \R$.
\begin{enumerate}[(i)]
\item
$\exp(xy) = {\bigl(\exp(x)\bigr)}^y$.
\item
If $x > 0$, then $\ln(x^y) = y \ln (x)$.
\end{enumerate}
\end{prop}

\begin{remark}
There are other equivalent ways to define the exponential and the logarithm.
A common way is to define $E$ as the solution to the differential equation
$E'(x) = E(x)$, $E(0) = 1$.  See \exampleref{example:picardexponential},
for a sketch of that approach.  Yet another
approach is to define the exponential function by
power series, see \exampleref{example:exponentialbypowerseries}.
\end{remark}

\begin{remark}
We proved the uniqueness of the functions $L$ and $E$ from
just the properties $L(1)=0$, $L'(x) = \nicefrac{1}{x}$
and the equivalent condition for the exponential
$E'(x) = E(x)$, $E(0) = 1$.  Existence also follows
from just these properties.
Alternatively, uniqueness also follows
from the laws of exponents, see the exercises.
\end{remark}

\subsection{Exercícios}

\begin{exercise}
Given a real number $y$ and $b > 0$, define $f \colon (0,\infty) \to \R$
and $g \colon \R \to \R$ as $f(x) \coloneqq x^y$ and $g(x) \coloneqq b^x$.  Show that $f$
and $g$ are differentiable and find their derivative.
\end{exercise}

\begin{samepage}
\begin{exercise} \label{exercise:otherlogbases}
Let $b > 0$, $b\neq 1$ be given.
\begin{enumerate}[a)]
\item
Show that for every $y > 0$, there exists a unique number $x$
such that $y = b^x$.  Define
the \emph{\myindex{logarithm base $b$}},
$\log_b \colon (0,\infty) \to \R$, by
$\log_b(y) \coloneqq x$.
\item
Show that $\log_b(x) = \frac{\ln(x)}{\ln(b)}$.
\item
Prove that if $c > 0$, $c \neq 1$, then
$\log_b(x) = \frac{\log_c(x)}{\log_c(b)}$.
\item
Prove $\log_b(xy) =
\log_b(x)+\log_b(y)$, and $\log_b(x^y) = y \log_b(x)$.
\end{enumerate}
\end{exercise}
\end{samepage}

\begin{exercise}[requires \sectionref{sec:taylor}]
Use \hyperref[thm:taylor]{Taylor's theorem} to study the remainder term and show that for
all $x \in \R$
\begin{equation*}
e^x = \sum_{n=0}^\infty \frac{x^n}{n!} .
\end{equation*}
Hint: Do not differentiate the series term by term (unless you prove that
you can do that).
\end{exercise}

\begin{exercise}
Use the geometric sum formula to show (for $t \neq -1$)
\begin{equation*}
1-t+t^2-\cdots+{(-1)}^n t^n = \frac{1}{1+t} - \frac{{(-1)}^{n+1}t^{n+1}}{1+t}.
\end{equation*}
Using this fact, show
\begin{equation*}
\ln (1+x) = \sum_{n=1}^\infty \frac{{(-1)}^{n+1}x^n}{n} 
\end{equation*}
for all $x \in (-1,1]$ (note that $x=1$ is included).  Finally,
find the limit of the alternating harmonic series
\begin{equation*}
\sum_{n=1}^\infty \frac{{(-1)}^{n+1}}{n} = 1 - \frac{1}{2} +
\frac{1}{3} - \frac{1}{4} + \cdots
\end{equation*}
\end{exercise}

\begin{exercise}
Show 
\begin{equation*}
e^x = \lim_{n\to\infty} {\left( 1 + \frac{x}{n} \right)}^n .
\end{equation*}
Hint: Take the logarithm.\\
Note: The expression 
${\left( 1 + \frac{x}{n} \right)}^n$ arises in compound interest
calculations.  It is the amount of money in a bank account after 1 year
if 1 dollar was deposited initially at interest rate $x$ (e.g., $x=0.01$ for
$1\%$) and the interest was compounded $n$
times during the year.  The exponential $e^x$ is the result of continuous
compounding.
\end{exercise}

\begin{samepage}
\begin{exercise}
\leavevmode
\begin{enumerate}[a)]
\item
Prove that for all integers $n \geq 2$,
\begin{equation*}
\sum_{k=2}^{n}
\frac{1}{k}
\leq
\ln (n)
\leq
\sum_{k=1}^{n-1}
\frac{1}{k} .
\end{equation*}
\item
Prove that the limit
\begin{equation*}
\gamma \coloneqq \lim_{n\to\infty}
\left( \left( \sum_{k=1}^{n} \frac{1}{k} \right) - \ln (n) \right)
\end{equation*}
exists.  This constant $\gamma$ is known as the
\emph{\myindex{Euler--Mascheroni constant}}%
\footnote{Named for the Swiss mathematician
\href{https://en.wikipedia.org/wiki/Leonhard_Euler}{Leonhard Euler}
(1707--1783)
and the Italian mathematician
\href{https://en.wikipedia.org/wiki/Lorenzo_Mascheroni}{Lorenzo Mascheroni}
(1750--1800).}.  It is not known if $\gamma$ is rational or not.
Approximately, $\gamma \approx 0.5772$.
\end{enumerate}
\end{exercise}
\end{samepage}

\begin{exercise}
Show
\begin{equation*}
\lim_{x\to\infty} \frac{\ln(x)}{x} = 0 .
\end{equation*}
\end{exercise}

\begin{exercise}
Show that $e^x$ is \emph{\myindex{convex}}, in other words, show that 
if $a \leq x \leq b$, then
$e^x \leq e^a \frac{b-x}{b-a} + e^b \frac{x-a}{b-a}$.
\end{exercise}

\begin{exercise}
Using the logarithm find
\begin{equation*}
%\lim_{n\to\infty} {\left( 1 + \nicefrac{1}{n} \right)}^n = e .
\lim_{n\to\infty} n^{1/n} .
\end{equation*}
\end{exercise}

\begin{exercise}
Show that $E(x) = e^x$ is the unique continuous function such that
$E(x+y) = E(x)E(y)$ and $E(1) = e$.   Similarly, prove that $L(x) = \ln(x)$
is the unique continuous
function defined on positive $x$ such that $L(xy) = L(x)+L(y)$
and $L(e) = 1$.
\end{exercise}

\begin{exercise}[requires \sectionref{sec:taylor}]\label{exercise:nonanalytic}
Since $(e^x)' = e^x$, it is easy to see that $e^x$ is
\myindex{infinitely differentiable}\index{differentiable!infinitely}
(has derivatives of all orders).  Define the function $f \colon \R \to \R$.
\begin{equation*}
f(x) \coloneqq \begin{cases}
e^{-1/x} & \text{se } x > 0, \\
0 & \text{se } x \leq 0.
\end{cases}
\end{equation*}
\begin{enumerate}[a)]
\item
Prove that for every $m \in \N$,
\begin{equation*}
\lim_{x \to 0^+} \frac{e^{-1/x}}{x^m} = 0 .
\end{equation*}
\item
Prove that $f$ is infinitely differentiable.
\item
Compute the Taylor series for $f$ at the origin, that is,
\begin{equation*}
\sum_{k=0}^\infty
\frac{f^{(k)}(0)}{k!}x^k .
\end{equation*}
Show that it converges, but show that it does not converge to $f(x)$
for any given $x > 0$.
\end{enumerate}
\end{exercise}

%%%%%%%%%%%%%%%%%%%%%%%%%%%%%%%%%%%%%%%%%%%%%%%%%%%%%%%%%%%%%%%%%%%%%%%%%%%%%%

\sectionnewpage
\section{Improper integrals}
\label{sec:impropriemann}

%mbxINTROSUBSECTION

\sectionnotes{2--3 lectures (optional section, can safely be skipped, 
requires the optional \sectionref{sec:limitatinf})}

Often it is necessary to integrate over the
entire real line, or an unbounded interval of the form $[a,\infty)$ or
$(-\infty,b]$.  We may also wish to integrate unbounded functions
defined on an open bounded interval $(a,b)$.
For such intervals or functions, the Riemann integral is not defined,
but we will write down
the integral anyway in the spirit of \lemmaref{lemma:boundedimpriemann}.
These integrals are called \emph{\myindex{improper integrals}}
and are limits
of integrals rather than integrals themselves.

\begin{defn}
Suppose $f \colon [a,b) \to \R$ is a function (not necessarily bounded)
that is Riemann integrable on $[a,c]$ for all $c < b$.  We define
\begin{equation*}
\int_a^b f \coloneqq \lim_{c \to b^-} \int_a^{c} f
\end{equation*}
if the limit exists.

Suppose $f \colon [a,\infty) \to \R$ is a function such that
$f$ is Riemann integrable on $[a,c]$ for all $c < \infty$.  
We define
\begin{equation*}
\int_a^\infty f \coloneqq \lim_{c \to \infty} \int_a^c f
\end{equation*}
if the limit exists.

If the limit exists, we say the improper integral
\emph{converges}\index{convergent!improper integral}.
If the limit does not exist, we say the improper integral
\emph{diverges}\index{divergent!improper integral}.

We similarly define improper integrals for the left-hand endpoint.
We leave this definition to the reader.
\end{defn}

For a finite endpoint $b$, if $f$ is bounded,
then
\lemmaref{lemma:boundedimpriemann} says that we defined nothing new.  What is new is that
we can apply this definition to unbounded functions.
The following set of examples is
so useful that we state it as a proposition.

\begin{prop}[$p$-test for integrals]%
\index{p-test for integrals@$p$-test for integrals}
\label{impropriemann:ptest}
The improper integral
\begin{equation*}
\int_1^\infty \frac{1}{x^p} \,dx
\end{equation*}
converges to $\frac{1}{p-1}$ if $p > 1$ and diverges if $0 < p \leq 1$.

The improper integral
\begin{equation*}
\int_0^1 \frac{1}{x^p} \,dx
\end{equation*}
converges to $\frac{1}{1-p}$ if $0 < p < 1$ and diverges if $p \geq 1$.
\end{prop}

\begin{proof}
The proof follows by application of the
\hyperref[thm:FTCv1]{fundamental theorem of calculus}.
Let us do the proof for $p > 1$ for the infinite right endpoint and
leave the rest to the reader.  Hint: You should handle $p=1$
separately.

Suppose $p > 1$.  Then
using the fundamental theorem,
\begin{equation*}
\int_1^b \frac{1}{x^p} \,dx
=
\int_1^b x^{-p} \,dx
=
\frac{b^{-p+1}}{-p+1}
-
\frac{1^{-p+1}}{-p+1}
=
\frac{-1}{(p-1)b^{p-1}}
+
\frac{1}{p-1} .
\end{equation*}
As $p > 1$, we have $p-1 > 0$.  Take the limit as $b \to \infty$
to obtain that $\frac{1}{b^{p-1}}$ goes to 0.  The result follows.
\end{proof}

We state the following proposition on \myquote{tails} for just one type
of improper integral, though the proof is straightforward
and the same for other types of improper integrals.

\begin{prop} \label{impropriemann:tail}
Let $f \colon [a,\infty) \to \R$ be a function
that is Riemann integrable on $[a,b]$ for all $b > a$.
For every $b > a$, the integral
$\int_b^\infty f$ converges if and only if $\int_a^\infty f$
converges, in which case
\begin{equation*}
\int_a^\infty f
=
\int_a^b f +
\int_b^\infty f .
\end{equation*}
\end{prop}

\begin{proof}
Let $c > b$.  Then
\begin{equation*}
\int_a^c f
=
\int_a^b f +
\int_b^c f .
\end{equation*}
Taking the limit $c \to \infty$ finishes the proof.
\end{proof}

Nonnegative functions are easier to work with
as the following proposition demonstrates.
The exercises will show that this proposition
holds only for nonnegative functions.
Analogues of this proposition
exist for all the other types of improper integrals and are left to the
student.

\begin{prop} \label{impropriemann:possimp}
Suppose $f \colon [a,\infty) \to \R$ is nonnegative ($f(x)
\geq 0$ for all $x$) and
$f$ is Riemann integrable on $[a,b]$ for all $b > a$.
\begin{enumerate}[(i)]
\item 
\begin{equation*}
\int_a^\infty f = \sup \left\{ \int_a^x f : x \geq a \right\} .
\end{equation*}
\item
Suppose $\{ x_n \}_{n=1}^\infty$
is a sequence with $\lim_{n\to\infty} x_n = \infty$.  Then
$\int_a^\infty f$ converges if and only if $\lim_{n\to\infty} \int_a^{x_n} f$ exists, in
which case
\begin{equation*}
\int_a^\infty f = \lim_{n\to\infty} \int_a^{x_n} f .
\end{equation*}
\end{enumerate}
\end{prop}

In the first item we allow for the value of $\infty$ in the
supremum indicating that the integral diverges to infinity.

\begin{proof}
We start with the first item.
As $f$ is nonnegative,
$\int_a^x f$ is increasing as a function of $x$.
If the supremum is infinite, then for every $M \in \R$
we find $N$ such that $\int_a^N f \geq M$.  As $\int_a^x f$
is increasing, $\int_a^x f \geq M$ for all $x \geq N$.  So
$\int_a^\infty f$ diverges to infinity.

Next suppose the supremum is finite, say
$A \coloneqq \sup \left\{ \int_a^x f : x \geq a \right\}$.
For every $\epsilon > 0$, we find an $N$ such that
$A - \int_a^N f < \epsilon$.  As $\int_a^x f$ is increasing,
then
$A - \int_a^x f < \epsilon$ for all $x \geq N$ and hence
$\int_a^\infty f$ converges to $A$.

Let us look at the second item.
If $\int_a^\infty f$ converges, then every sequence $\{ x_n \}_{n=1}^\infty$ going to
infinity works.  The trick is
proving the other direction.  Suppose $\{ x_n \}_{n=1}^\infty$ is such that
$\lim_{n\to\infty} x_n = \infty$ and
\begin{equation*}
\lim_{n\to\infty} \int_a^{x_n} f = A
\end{equation*}
converges.  Given $\epsilon > 0$, pick $N$ such that for
all $n \geq N$, we have
$A - \epsilon < \int_a^{x_n} f < A + \epsilon$.
Because $\int_a^x f$ is increasing as a function of $x$, we have that for all
$x \geq x_N$
\begin{equation*}
A - \epsilon < \int_a^{x_N} f \leq \int_a^x f .
\end{equation*}
As $\{ x_n \}_{n=1}^\infty$ goes to $\infty$, we have that for any 
$x$, there is an $x_m$ such that $m \geq N$ and $x \leq x_m$.  Then
\begin{equation*}
\int_a^{x} f \leq \int_a^{x_m} f < A + \epsilon .
\end{equation*}
In particular, for all $x \geq x_N$, we have
$\abs{\int_a^{x} f - A} < \epsilon$.
\end{proof}

\begin{prop}[Comparison test for improper integrals]%
\index{comparison test for improper integrals}
Let
$f \colon [a,\infty) \to \R$ and
$g \colon [a,\infty) \to \R$ be functions
that are Riemann integrable on $[a,b]$ for all $b > a$.   Suppose
that for all $x \geq a$,
\begin{equation*}
\babs{f(x)} \leq g(x) .
\end{equation*}
\begin{enumerate}[(i)]
\item If $\int_a^\infty g$ converges, then $\int_a^\infty f$ converges,
and in this case 
$\abs{\int_a^\infty f} \leq \int_a^\infty g$.
\item If $\int_a^\infty f$ diverges, then $\int_a^\infty g$ diverges.
\end{enumerate}
\end{prop}

\begin{proof}
We start with the first item.
For every $b$ and $c$, such that $a \leq b \leq c$, we have 
$-g(x) \leq f(x) \leq g(x)$, and so
\begin{equation*}
\int_b^c -g \leq \int_b^c f \leq \int_b^c g  .
\end{equation*}
In other words, $\abs{\int_b^c f} \leq \int_b^c g$.

Let $\epsilon > 0$ be given.  Because
of \propref{impropriemann:tail},
\begin{equation*}
\int_a^\infty g =
\int_a^b g +
\int_b^\infty g .
\end{equation*}
As $\int_a^b g$ goes to
$\int_a^\infty g$ as $b$ goes to infinity,
$\int_b^\infty g$ goes to 0 as $b$ goes to infinity.  Choose $B$
such that
\begin{equation*}
\int_B^\infty g < \epsilon .
\end{equation*}
As $g$ is nonnegative, if $B \leq b < c$, then
$\int_b^c g < \epsilon$ as well.
Let $\{ x_n \}_{n=1}^\infty$ be a sequence going to infinity.  Let $M$ be such that
$x_n \geq B$ for all $n \geq M$.  Take $n, m \geq M$,
with $x_n \leq x_m$,
\begin{equation*}
\abs{\int_a^{x_m} f - \int_a^{x_n} f} 
=
\abs{\int_{x_n}^{x_m} f} 
\leq \int_{x_n}^{x_m} g < \epsilon .
\end{equation*}
Therefore, the sequence $\bigl\{ \int_a^{x_n} f \bigr\}_{n=1}^\infty$ is Cauchy and hence converges.

We need to show that the limit is unique.  Suppose $\{ x_n \}_{n=1}^\infty$ is a sequence
converging to infinity such that
$\bigl\{ \int_a^{x_n} f \bigr\}_{n=1}^\infty$ converges to $L_1$, and $\{
y_n \}_{n=1}^\infty$ is a sequence
converging to infinity such that
$\bigl\{ \int_a^{y_n} f \bigr\}_{n=1}^\infty$ converges to $L_2$.  Then there must be some $n$ such
that
$\babs{\int_a^{x_n} f - L_1} < \epsilon$ and 
$\babs{\int_a^{y_n} f - L_2} < \epsilon$.  We can also suppose $x_n \geq B$
and $y_n \geq B$.  Then
\begin{equation*}
\sabs{L_1 - L_2} \leq
\abs{L_1 - \int_a^{x_n} f}
+
\abs{\int_a^{x_n} f- \int_a^{y_n} f}
+
\abs{\int_a^{y_n} f - L_2}
<
\epsilon
+
\abs{\int_{x_n}^{y_n} f}
+
\epsilon
<
3 \epsilon.
\end{equation*}
As $\epsilon > 0$ was arbitrary, $L_1 = L_2$, and hence
$\int_a^\infty f$ converges.
Above we have shown that $\abs{\int_a^c f} \leq \int_a^c g$ for all $c > a$.
By taking the limit $c \to \infty$, the first item is proved.

The second item is simply a contrapositive of the first item.
\end{proof}

\begin{example}
The improper integral
\begin{equation*}
\int_0^\infty \frac{\sin(x^2)(x+2)}{x^3+1} \,dx
\end{equation*}
converges.

Proof:  Observe we simply need to show
that the integral converges when going from 1 to infinity.
For $x \geq 1$ we obtain
\begin{equation*}
\abs{\frac{\sin(x^2)(x+2)}{x^3+1}}
\leq
\frac{x+2}{x^3+1}
\leq \frac{x+2}{x^3} \leq
\frac{x+2x}{x^3} \leq \frac{3}{x^2} .
\end{equation*}
Then
\begin{equation*}
\int_1^\infty \frac{3}{x^2}\,dx
=
3 \int_1^\infty \frac{1}{x^2}\,dx
%=
%\lim_{c\to\infty} \int_1^c \frac{3}{x^2} \,dx
=
3 .
\end{equation*}
So using the comparison test and the tail test, the original
integral converges.
\end{example}

\begin{example}
You should be careful when doing formal manipulations with improper
integrals.
The integral
\begin{equation*}
\int_2^\infty \frac{2}{x^2-1}\,dx
\end{equation*}
converges via the comparison test using $\nicefrac{1}{x^2}$ again.  However, if you
succumb to the temptation to write
\begin{equation*}
\frac{2}{x^2-1} = 
\frac{1}{x-1}
-
\frac{1}{x+1} 
\end{equation*}
and try to integrate each part separately, you will not succeed.
It is \emph{not} true that you can split the improper
integral in two; you cannot split the limit.
\begin{equation*}
\begin{split}
\int_2^\infty \frac{2}{x^2-1} \,dx &=
\lim_{b\to \infty} \int_2^b \frac{2}{x^2-1} \,dx
\\
&=
\lim_{b\to \infty}
\left(
\int_2^b \frac{1}{x-1}\,dx
-
\int_2^b \frac{1}{x+1}\,dx
\right)
\\
&\neq
\int_2^\infty \frac{1}{x-1}\,dx
-
\int_2^\infty \frac{1}{x+1}\,dx .
\end{split}
\end{equation*}
The last line in the computation does not even make sense.  Both of the
integrals diverge to infinity, since we can
apply the comparison test appropriately with
$\nicefrac{1}{x}$.  We get $\infty - \infty$.
\end{example}

Now suppose we need to take limits at both endpoints.

\begin{defn}
Suppose $f \colon (a,b) \to \R$ is a function
that is Riemann integrable on $[c,d]$ for all $c$, $d$
such that $a < c < d < b$, then we define
\begin{equation*}
\int_a^b f \coloneqq \lim_{c \to a^+} \, \lim_{d \to b^-} \, \int_{c}^{d} f
\end{equation*}
if the limits exist.

Suppose $f \colon \R \to \R$ is a function such that
$f$ is Riemann integrable on all bounded intervals $[a,b]$.  Then
we define
\begin{equation*}
\int_{-\infty}^\infty f \coloneqq \lim_{c \to -\infty} \, \lim_{d \to \infty} \, \int_c^d f
\end{equation*}
if the limits exist.

We similarly define improper integrals with one infinite and one finite
improper endpoint.  We leave this definition to the reader.
\end{defn}

One ought to always be careful about double limits.  The definition
given above says that we first take the limit as $d$ goes to $b$ or
$\infty$ for a fixed $c$, and then we take the limit in $c$.
We will have to prove that in this case it does not matter which limit
we compute first.

\begin{example}
\begin{equation*}
\int_{-\infty}^\infty \frac{1}{1+x^2} \, dx
=
\lim_{a \to -\infty} \, \lim_{b \to \infty} \,
\int_{a}^b \frac{1}{1+x^2} \, dx
=
\lim_{a \to -\infty} \, \lim_{b \to \infty}
\bigl( \arctan(b) - \arctan(a) \bigr)
=
\pi .
\end{equation*}
\end{example}

In the definition, the order of the limits can always be switched if they
exist.  Let us state and prove this fact only for the limits at infinity.

\begin{prop}
Suppose $f \colon \R \to \R$ is integrable on every bounded interval
$[a,b]$.
Then 
\begin{equation*}
\lim_{a \to -\infty} \, \lim_{b \to \infty} \, \int_a^b f
\quad \text{converges} \qquad \text{if and only if} \qquad
\lim_{b \to \infty}
\,
\lim_{a \to -\infty}
\,
\int_a^b f
\quad
\text{converges,}
\end{equation*}
in which case the two
expressions are equal.  If either of the
expressions converges, then the improper integral converges and
\begin{equation*}
\lim_{a\to\infty}
\int_{-a}^a f
=
\int_{-\infty}^\infty f .
\end{equation*}
\end{prop}

\begin{proof}
Without loss of generality, assume $a < 0$ and $b > 0$.  Suppose
the first expression converges.  Then
\begin{equation*}
\begin{split}
\lim_{a \to -\infty} \, \lim_{b \to \infty} \, \int_a^b f
& =
\lim_{a \to -\infty} \, \lim_{b \to \infty}
\left(
\int_a^0 f
+
\int_0^b f
\right)
=
\left(
\lim_{a \to -\infty}
\int_a^0 f
\right)
+
\left(
 \lim_{b \to \infty}
\int_0^b f
\right) \\
& = 
 \lim_{b \to \infty}
\left(
\left(
\lim_{a \to -\infty}
\int_a^0 f
\right) 
+
\int_0^b f
\right)
=
 \lim_{b \to \infty} \,
\lim_{a \to -\infty}
\left(
\int_a^0 f
+
\int_0^b f
\right)  .
\end{split}
\end{equation*}
Similar computation shows the other direction.  Therefore, if
either expression converges, then the improper integral converges
and
\begin{multline*}
\int_{-\infty}^\infty f
=
\lim_{a \to -\infty} \, \lim_{b \to \infty} \, \int_a^b f
=
\left(
\lim_{a \to -\infty}
\int_a^0 f
\right)
+
\left(
 \lim_{b \to \infty}
\int_0^b f
\right)
\\
=
\left(
\lim_{a \to \infty}
\int_{-a}^0 f
\right)
+
\left(
 \lim_{a \to \infty}
\int_0^a f
\right)
=
\lim_{a \to \infty}
\left(
\int_{-a}^0 f
+
\int_0^a f
\right)
=
\lim_{a \to \infty}
\int_{-a}^a f .
\end{multline*}
\end{proof}

\begin{example}
On the other hand, you must be careful to
take the limits independently before you know convergence.  Let
$f(x) = \frac{x}{\sabs{x}}$ for $x \neq 0$ and $f(0) = 0$.
If $a < 0$ and $b > 0$, then
\begin{equation*}
\int_{a}^b f
=
\int_{a}^0 f
+
\int_{0}^b f
=
a+b .
\end{equation*}
For every fixed $a < 0$, the limit as $b \to \infty$ is infinite.
So even the first limit does not exist,
and the improper integral $\int_{-\infty}^\infty f$ does not converge.
On the other hand, if $a > 0$, then
\begin{equation*}
\int_{-a}^{a} f
=
(-a)+a = 0 .
\end{equation*}
Therefore,
\begin{equation*}
\lim_{a\to\infty}
\int_{-a}^{a} f
= 0 .
\end{equation*}
\end{example}

\begin{example}
An example to keep in mind for improper integrals
is the \emph{\myindex{sinc function}}%
\footnote{Shortened from Latin: \emph{sinus cardinalis}}.
This function comes up quite often
in both pure and applied mathematics.  Define
\begin{equation*}
\operatorname{sinc}(x) \coloneqq
\begin{cases}
\frac{\sin(x)}{x} & \text{se } x \neq 0 , \\
1 & \text{se } x = 0 .
\end{cases}
\end{equation*}
\begin{myfigureht}
\myincludegraphics{sincfig}{%
A graph of the sinc function on the interval from minus 4 pi to 4 pi.
It oscillates sort of like the sine function to the right
of the origin, and sort of like minus sine to the left of the
origin.  The root at the origin disappears and the function is simply 1
at the origin.  The oscillations quickly get smaller and smaller
as we go away from the origin.  The function is clearly even.}
\caption{The sinc function.\label{figsinc}}
\end{myfigureht}

It is not difficult to show that
the sinc function is continuous at zero, but that is
not important right now.  What is important is that
\begin{equation*}
\int_{-\infty}^\infty \operatorname{sinc}(x) \,dx = \pi ,
\qquad \text{while} \qquad
\int_{-\infty}^\infty \babs{\operatorname{sinc}(x)} \,dx = \infty .
\end{equation*}
The integral of the sinc function is a continuous analogue of the
alternating harmonic series $\sum_{n=1}^\infty \nicefrac{{(-1)}^n}{n}$, while the
absolute value is like the regular harmonic series $\sum_{n=1}^\infty \nicefrac{1}{n}$.
In particular, the fact that the integral converges must be done directly
rather than using the comparison test.

We will not prove the first statement exactly.  Let us simply prove
that the integral of the sinc function converges, but we will not worry
about the exact limit.  Because $\frac{\sin(-x)}{-x} = \frac{\sin(x)}{x}$, it is
enough to show that
\begin{equation*}
\int_{2\pi}^\infty \frac{\sin(x)}{x}\,dx
\end{equation*}
converges.  We
also avoid $x=0$ this way to make our life simpler.

For every $n \in \N$, we have that for $x \in [\pi 2n, \pi (2n+1)]$,
\begin{equation*}
\frac{\sin(x)}{\pi (2n+1)}
\leq
\frac{\sin(x)}{x}
\leq
\frac{\sin(x)}{\pi 2n} ,
\end{equation*}
as $\sin(x) \geq 0$.  For $x \in [\pi (2n+1), \pi (2n+2)]$,
\begin{equation*}
\frac{\sin(x)}{\pi (2n+1)}
\leq
\frac{\sin(x)}{x}
\leq
\frac{\sin(x)}{\pi (2n+2)} ,
\end{equation*}
as $\sin(x) \leq 0$.

Via the fundamental theorem of calculus,
\begin{equation*}
%mbxlatex \begin{aligned}
\frac{2}{\pi (2n+1)}
=
\int_{\pi 2n}^{\pi (2n+1)}
\frac{\sin(x)}{\pi (2n+1)}
\,dx
%mbxlatex &
\leq
\int_{\pi 2n}^{\pi (2n+1)}
\frac{\sin(x)}{x}
\,dx
%mbxlatex \\
%mbxlatex &
\leq
\int_{\pi 2n}^{\pi (2n+1)}
\frac{\sin(x)}{\pi 2n}
\,dx
=
\frac{1}{\pi n} .
%mbxlatex \end{aligned}
\end{equation*}
Similarly,
\begin{equation*}
\frac{-2}{\pi (2n+1)}
\leq
\int_{\pi (2n+1)}^{\pi (2n+2)}
\frac{\sin(x)}{x}
\,dx
\leq
\frac{-1}{\pi (n+1)} .
\end{equation*}
Adding the two together we find
\begin{equation*}
%mbxlatex \begin{aligned}
0
=
\frac{2}{\pi (2n+1)}
+
\frac{-2}{\pi (2n+1)}
%mbxlatex &
\leq
\int_{2\pi n}^{2\pi (n+1)}
\frac{\sin(x)}{x}
\,dx
%mbxlatex \\
%mbxlatex &
\leq
\frac{1}{\pi n} 
+
\frac{-1}{\pi (n+1)} 
=
\frac{1}{\pi n(n+1)} .
%mbxlatex \end{aligned}
\end{equation*}
See \figureref{fig:sincbound}.
\begin{myfigureht}
\myincludegraphics{sincbound}{%
A graph of several functions on the interval
from pi times 2 n to pi times quantity 2 n plus 2
is shown.  They all look like a sine wave that is losing amplitude.
A dark bold line is the function sine of x all over x
graphed over the entire interval.
Like all the other functions it is zero at the left-hand endpoint,
the right-hand endpoint and the middle point pi times quantity 2 n
plus 1.  A dotted line on the first subinterval (the left half) 
very close to and below the solid line denotes the
function sine of x all over quantity pi times quantity 2 n plus 1.
A dashed line over the first subinterval very close to and above
the solid line denotes the function sine of x over quantity pi 2 n.
The area between the dashed line and the x-axis over the first
subinterval is shaded.
On the right half of the interval, a dotted line denotes
the same function as before, and again it is below the solid line.
A dash-dotted line that is slightly above the solid line
now denotes the function sine of x all over quantity pi times
quantity 2 n plus 2.  The area above the dash-dotted line is shaded.}
\caption{Bound of $\int_{2\pi n}^{2\pi (n+1)} \frac{\sin(x)}{x} \,dx$ using
the shaded integral (signed area
$\frac{1}{\pi n} 
+
\frac{-1}{\pi (n+1)}$).\label{fig:sincbound}}
\end{myfigureht}

For $k \in \N$, 
\begin{equation*}
\int_{2\pi}^{2k\pi} \frac{\sin(x)}{x} \,dx
=
\sum_{n=1}^{k-1}
\int_{2\pi n}^{2\pi (n+1)} \frac{\sin(x)}{x} \,dx 
\leq
\sum_{n=1}^{k-1}
\frac{1}{\pi n(n+1)} .
\end{equation*}
We find the partial sums of a series with positive terms.
The series
converges as
$\sum_{n=1}^\infty \frac{1}{\pi n (n+1)}$ is a convergent series.  Thus
as a sequence,
\begin{equation*}
\lim_{k\to \infty} \int_{2\pi}^{2k\pi} \frac{\sin(x)}{x} \,dx
=L \leq
\sum_{n=1}^{\infty}
\frac{1}{\pi n(n+1)} < \infty .
\end{equation*}

Let $M > 2\pi$ be arbitrary, and let $k \in \N$
be the largest integer such that $2k\pi \leq M$.
For $x \in [2k\pi,M]$, we have 
$\frac{-1}{2k\pi} \leq \frac{\sin(x)}{x} \leq \frac{1}{2k\pi}$, and so
\begin{equation*}
\abs{\int_{2k\pi}^{M} \frac{\sin(x)}{x} \,dx }  \leq
\frac{M-2k\pi}{2k\pi} \leq \frac{1}{k} .
\end{equation*}
As $k$ is the largest $k$ such that $2k\pi \leq M$,
then as $M\in \R$ goes to infinity, so does $k \in \N$.

Then
\begin{equation*}
\int_{2\pi}^M \frac{\sin(x)}{x}\,dx
=
\int_{2\pi}^{2k\pi} \frac{\sin(x)}{x} \,dx
+
\int_{2k\pi}^{M} \frac{\sin(x)}{x} \,dx .
\end{equation*}
As $M$ goes to infinity,
the first term on the
right-hand side goes to $L$,
and the second term on the
right-hand side
goes to zero.  Hence,
\begin{equation*}
\int_{2\pi}^\infty \frac{\sin(x)}{x} \,dx = L .
%\leq \sum_{n=1}^{\infty}
%\frac{1}{\pi n(n+1)} < \infty .
\end{equation*}

The double-sided integral of sinc also exists as noted above.
We leave the other statement---that the integral
of the absolute value of the sinc function diverges---as an exercise.
\end{example}

\subsection{Integral test for series}

The fundamental theorem 
of calculus can be used in proving a series is summable and to estimate its sum.

\begin{prop}[Integral test]\index{integral test for series}
Suppose $f \colon [k,\infty) \to \R$ is a decreasing nonnegative
function where $k \in \Z$.  Then
\begin{equation*}
\sum_{n=k}^\infty f(n)
\quad \text{converges} \qquad \text{if and only if}
\qquad
\int_k^\infty f
\quad \text{converges}.
\end{equation*}
In this case 
\begin{equation*}
\int_k^\infty f
\leq
\sum_{n=k}^\infty f(n)
\leq
f(k)+
\int_k^\infty f .
\end{equation*}
\end{prop}
See \figureref{fig:integraltest}, for an illustration with $k=1$.
By \propref{prop:monotoneintegrable},
$f$ is integrable
on every interval $[k,b]$ for all $b > k$, so the statement of the theorem
makes sense without additional hypotheses of integrability.
\begin{myfigureht}
\myincludegraphics{integraltest}{%
A diagram of the Darboux rectangles for a decreasing
function starting at x equals 1 and going in intervals
of 1 until x equals 10 on the right.  As the function
is decreasing, we notice that the function graph goes through
the top left corner of the entire rectangle and also the top
right corner of the shaded rectangle.}
\caption{The area under the curve, 
$\int_1^\infty f$, is bounded below
by the area of the shaded rectangles,
$f(2)+f(3)+f(4)+\cdots$, and bounded above
by the area of the entire rectangles, both the shaded and unshaded parts,
$f(1)+f(2)+f(3)+\cdots$.\label{fig:integraltest}}
\end{myfigureht}

\begin{proof}
Let $\ell, m \in \Z$ be such that $m > \ell \geq k$.
Because $f$ is decreasing, we have
$\int_{n}^{n+1} f \leq f(n) \leq \int_{n-1}^{n} f$.  Therefore,
\begin{equation} \label{impropriemann:eqseries}
\int_\ell^m f
=
\sum_{n=\ell}^{m-1} \int_{n}^{n+1} f
\leq
\sum_{n=\ell}^{m-1} f(n)
\leq
f(\ell) +
\sum_{n=\ell+1}^{m-1} \int_{n-1}^{n} f
\leq
f(\ell)+
\int_\ell^{m-1} f .
\end{equation}

Suppose first that $\int_k^\infty f$ converges and
let $\epsilon > 0$ be given.
As before, since $f$ is positive, there exists
an $L \in \N$ such that if $\ell \geq L$, then
$\int_\ell^{m} f < \nicefrac{\epsilon}{2}$ for all $m \geq \ell$.
The function 
$f$ must decrease to zero (why?), so make $L$ large enough so that
for $\ell \geq L$, we have $f(\ell) < \nicefrac{\epsilon}{2}$.
Thus, for $m > \ell \geq L$, we have
via \eqref{impropriemann:eqseries},
\begin{equation*}
\sum_{n=\ell}^{m} f(n)
\leq
f(\ell)+
\int_\ell^{m} f < \nicefrac{\epsilon}{2} + \nicefrac{\epsilon}{2} = \epsilon .
\end{equation*}
The series is therefore Cauchy and thus converges.  The estimate in the
proposition is obtained by letting $m$ go to infinity in
\eqref{impropriemann:eqseries} with $\ell = k$.

Conversely, suppose $\int_k^\infty f$ diverges.  
As $f$ is positive, by \propref{impropriemann:possimp},
the sequence $\{ \int_k^m f \}_{m=k}^\infty$ diverges to infinity.
Using
\eqref{impropriemann:eqseries} with $\ell = k$, we find
\begin{equation*}
\int_k^m f
\leq
\sum_{n=k}^{m-1} f(n) .
\end{equation*}
As the left-hand side goes to infinity as $m \to \infty$, so does the
right-hand side.
\end{proof}

\begin{example}
The integral test can be used not only to show that a series converges, but to
estimate its sum to arbitrary precision.
Let us show $\sum_{n=1}^\infty \frac{1}{n^2}$ exists and
estimate its sum to within 0.01.  As this series is the $p$-series for
$p=2$, we already proved it converges (let us pretend we do not know that),
but we only roughly estimated its sum.

The fundamental theorem of calculus says that for all $k \in \N$,
\begin{equation*}
\int_{k}^\infty \frac{1}{x^2}\,dx = \frac{1}{k} .
\end{equation*}
In particular, the series must converge.  But we also have
\begin{equation*}
\frac{1}{k} = \int_k^\infty \frac{1}{x^2}\,dx
\leq
\sum_{n=k}^\infty \frac{1}{n^2}
\leq
\frac{1}{k^2}
+
\int_k^\infty \frac{1}{x^2}\,dx
=
\frac{1}{k^2}
+
\frac{1}{k} .
\end{equation*}
Adding the partial sum up to $k-1$ we get
\begin{equation*}
\frac{1}{k} + \sum_{n=1}^{k-1} \frac{1}{n^2}
\leq
\sum_{n=1}^\infty \frac{1}{n^2}
\leq
\frac{1}{k^2}
+
\frac{1}{k} + \sum_{n=1}^{k-1} \frac{1}{n^2} .
\end{equation*}
In other words,
$\nicefrac{1}{k} + \sum_{n=1}^{k-1} \nicefrac{1}{n^2}$ is an estimate for
the sum to within $\nicefrac{1}{k^2}$.  Therefore, if we wish to
find the sum to within 0.01, we note $\nicefrac{1}{{10}^2} = 0.01$.  We
obtain
\begin{equation*}
1.6397\ldots
\approx
\frac{1}{10} + \sum_{n=1}^{9} \frac{1}{n^2}
\leq
\sum_{n=1}^\infty \frac{1}{n^2}
\leq
\frac{1}{100}
+
\frac{1}{10} + \sum_{n=1}^{9} \frac{1}{n^2}
\approx
1.6497\ldots .
\end{equation*}
The actual sum is $\nicefrac{\pi^2}{6} \approx 1.6449\ldots$. 
\end{example}

\subsection{Exercícios}

\begin{exercise}
Finish the proof of \propref{impropriemann:ptest}.
\end{exercise}

\begin{exercise}
Find out for which $a \in \R$ does $\sum_{n=1}^\infty e^{an}$ converge.
When the series converges, find an upper bound for the sum.
\end{exercise}

\begin{exercise}
\leavevmode
\begin{enumerate}[a)]
\item
Estimate $\sum_{n=1}^\infty \frac{1}{n(n+1)}$ correct to within 0.01
using the integral test.
\item
Compute the limit of the series exactly
and compare.  Hint: The sum telescopes.
\end{enumerate}
\end{exercise}

\begin{exercise}
Prove 
\begin{equation*}
\int_{-\infty}^\infty \babs{\operatorname{sinc}(x)}\,dx = \infty .
\end{equation*}
Hint: Again, it is enough to show this on just one side.
\end{exercise}

\begin{exercise}
Can you interpret
\begin{equation*}
\int_{-1}^1 \frac{1}{\sqrt{\sabs{x}}}\,dx
\end{equation*}
as an improper integral?  If so, compute its value.
\end{exercise}

\begin{exercise}
Take $f \colon [0,\infty) \to \R$, Riemann integrable on
every interval $[0,b]$, and such that there exist $M$, $a$, and $T$,
such that $\babs{f(t)} \leq M e^{at}$ for all $t \geq T$.  Show that the
\emph{\myindex{Laplace transform}} of $f$ exists.  That is, for
every $s > a$ the following integral converges:
\begin{equation*}
F(s) \coloneqq \int_{0}^\infty f(t) e^{-st} \,dt .
\end{equation*}
\end{exercise}

\begin{exercise}
Let $f \colon \R \to \R$ be a Riemann integrable function
on every interval $[a,b]$, and such
that $\int_{-\infty}^\infty \babs{f(x)}\,dx < \infty$.  Show that the
\emph{\myindex{Fourier sine and cosine transforms}}
exist.  That is, for every $\omega \geq 0$ the
following integrals converge
\begin{equation*}
F^s(\omega) \coloneqq \frac{1}{\pi} \int_{-\infty}^\infty f(t) \sin(\omega t) \,dt ,
\qquad
F^c(\omega) \coloneqq \frac{1}{\pi} \int_{-\infty}^\infty f(t) \cos(\omega t) \,dt .
\end{equation*}
Furthermore, show that $F^s$ and $F^c$ are bounded functions.
\end{exercise}

\begin{exercise}
Suppose $f \colon [0,\infty) \to \R$ is Riemann integrable on every interval
$[0,b]$.  Show that  $\int_0^\infty f$ converges if and only if
for every $\epsilon > 0$ there exists an $M$ such that if $M \leq a < b$,
then $\babs{\int_a^b f} < \epsilon$.
\end{exercise}

\begin{exercise}
Suppose $f \colon [0,\infty) \to \R$ is nonnegative and
\emph{decreasing}.  Prove:
\begin{enumerate}[a)]
\item
If $\int_0^\infty f < \infty$, then $\lim\limits_{x\to\infty} f(x) = 0$.
\item
The converse does not hold.
\end{enumerate}
\end{exercise}

\begin{exercise}
Find an example of an \emph{unbounded} continuous function $f \colon
[0,\infty) \to \R$ that is nonnegative and such that $\int_0^\infty f < \infty$.
Note that $\lim_{x\to\infty} f(x)$ will not exist; compare
previous exercise.
Hint: On each interval $[k,k+1]$, $k \in \N$, define a function whose
integral over this interval is less than say $2^{-k}$.
\end{exercise}

\begin{exercise}[More challenging]
Find an example of a function $f \colon [0,\infty) \to \R$ integrable on all
intervals such that $\lim_{n\to\infty} \int_0^n f$ exists as a
limit of a sequence (so $n \in \N$), but such that
$\int_0^\infty f$ does not exist.
Hint: For all $n\in \N$, divide $[n,n+1]$ into two halves.  On one half
make the function negative, on the other make the function positive.
\end{exercise}

\begin{exercise}
Suppose $f \colon [1,\infty) \to \R$ is such that
$g(x) \coloneqq x^2 f(x)$ is a bounded function. Prove that
$\int_1^\infty f$ converges.
\end{exercise}

\begin{exnote}
It is sometimes desirable to assign a value to integrals that normally
cannot be interpreted even as improper integrals,
e.g.\ $\int_{-1}^1 \nicefrac{1}{x}\,dx$.
Suppose $f \colon [a,b] \to \R$ is a function and $a < c < b$,
where $f$ is Riemann integrable on the intervals
$[a,c-\epsilon]$ and $[c+\epsilon,b]$ for all $\epsilon > 0$.
Define
the \emph{\myindex{Cauchy principal value}} of $\int_a^b f$ as
\begin{equation*}
p.v.\!\int_a^b f \coloneqq \lim_{\epsilon\to 0^+}
\left(
\int_a^{c-\epsilon} f + 
\int_{c+\epsilon}^b f
\right) ,
\end{equation*}
if the limit exists.
\end{exnote}

%\begin{samepage}
\begin{exercise}
\leavevmode
\begin{enumerate}[a)]
\item
Compute $p.v.\!\int_{-1}^1 \nicefrac{1}{x}\,dx$.
\item
Compute
$\lim_{\epsilon\to 0^+}
( \int_{-1}^{-\epsilon} \nicefrac{1}{x}\,dx + 
\int_{2\epsilon}^1 \nicefrac{1}{x}\,dx )$ and show it is not equal
to the principal value.
\item
Show that if $f$ is integrable on $[a,b]$, then
$p.v.\!\int_a^b f = \int_a^b f$ (for an arbitrary $c \in (a,b)$).
\item
Suppose $f \colon [-1,1] \to \R$
is an odd function ($f(-x)=-f(x)$) that is integrable on
$[-1,-\epsilon]$ and $[\epsilon,1]$ for all $\epsilon >0$.
Prove that 
$p.v.\!\int_{-1}^1 f = 0$
\item
Suppose 
$f \colon [-1,1] \to \R$ is continuous and differentiable at 0.  Show that
$p.v.\!\int_{-1}^1 \frac{f(x)}{x}\,dx$ exists.
\end{enumerate}
\end{exercise}
%\end{samepage}

\begin{samepage}
\begin{exercise}
Let $f \colon \R \to \R$ and 
$g \colon \R \to \R$ be continuous functions, where
$g(x) = 0$ for all $x \notin [a,b]$ for some interval $[a,b]$.
\begin{enumerate}[a)]
\item
Show that the
\emph{\myindex{convolution}}
\begin{equation*}
(g * f)(x) \coloneqq \int_{-\infty}^\infty f(t)g(x-t)\,dt 
\end{equation*}
is well-defined for all $x \in \R$.
\item
Suppose $\int_{-\infty}^\infty \babs{f(x)}\,dx < \infty$.  Prove that
\begin{equation*}
\lim_{x \to -\infty} (g * f)(x) = 0, \qquad \text{e} \qquad
\lim_{x \to \infty} (g * f)(x) = 0 .
\end{equation*}
\end{enumerate}
\end{exercise}
\end{samepage}
